% Options for packages loaded elsewhere
% Options for packages loaded elsewhere
\PassOptionsToPackage{unicode}{hyperref}
\PassOptionsToPackage{hyphens}{url}
\PassOptionsToPackage{dvipsnames,svgnames,x11names}{xcolor}
%
\documentclass[
  spanish,
  a4paper,
  oneside]{scrbook}
\usepackage{xcolor}
\usepackage[top=25mm,bottom=25mm,left=30mm,right=20mm]{geometry}
\usepackage{amsmath,amssymb}
\setcounter{secnumdepth}{5}
\usepackage{iftex}
\ifPDFTeX
  \usepackage[T1]{fontenc}
  \usepackage[utf8]{inputenc}
  \usepackage{textcomp} % provide euro and other symbols
\else % if luatex or xetex
  \usepackage{unicode-math} % this also loads fontspec
  \defaultfontfeatures{Scale=MatchLowercase}
  \defaultfontfeatures[\rmfamily]{Ligatures=TeX,Scale=1}
\fi
\usepackage{lmodern}
\ifPDFTeX\else
  % xetex/luatex font selection
\fi
% Use upquote if available, for straight quotes in verbatim environments
\IfFileExists{upquote.sty}{\usepackage{upquote}}{}
\IfFileExists{microtype.sty}{% use microtype if available
  \usepackage[]{microtype}
  \UseMicrotypeSet[protrusion]{basicmath} % disable protrusion for tt fonts
}{}
\makeatletter
\@ifundefined{KOMAClassName}{% if non-KOMA class
  \IfFileExists{parskip.sty}{%
    \usepackage{parskip}
  }{% else
    \setlength{\parindent}{0pt}
    \setlength{\parskip}{6pt plus 2pt minus 1pt}}
}{% if KOMA class
  \KOMAoptions{parskip=half}}
\makeatother
% Make \paragraph and \subparagraph free-standing
\makeatletter
\ifx\paragraph\undefined\else
  \let\oldparagraph\paragraph
  \renewcommand{\paragraph}{
    \@ifstar
      \xxxParagraphStar
      \xxxParagraphNoStar
  }
  \newcommand{\xxxParagraphStar}[1]{\oldparagraph*{#1}\mbox{}}
  \newcommand{\xxxParagraphNoStar}[1]{\oldparagraph{#1}\mbox{}}
\fi
\ifx\subparagraph\undefined\else
  \let\oldsubparagraph\subparagraph
  \renewcommand{\subparagraph}{
    \@ifstar
      \xxxSubParagraphStar
      \xxxSubParagraphNoStar
  }
  \newcommand{\xxxSubParagraphStar}[1]{\oldsubparagraph*{#1}\mbox{}}
  \newcommand{\xxxSubParagraphNoStar}[1]{\oldsubparagraph{#1}\mbox{}}
\fi
\makeatother


\usepackage{longtable,booktabs,array}
\usepackage{calc} % for calculating minipage widths
% Correct order of tables after \paragraph or \subparagraph
\usepackage{etoolbox}
\makeatletter
\patchcmd\longtable{\par}{\if@noskipsec\mbox{}\fi\par}{}{}
\makeatother
% Allow footnotes in longtable head/foot
\IfFileExists{footnotehyper.sty}{\usepackage{footnotehyper}}{\usepackage{footnote}}
\makesavenoteenv{longtable}
\usepackage{graphicx}
\makeatletter
\newsavebox\pandoc@box
\newcommand*\pandocbounded[1]{% scales image to fit in text height/width
  \sbox\pandoc@box{#1}%
  \Gscale@div\@tempa{\textheight}{\dimexpr\ht\pandoc@box+\dp\pandoc@box\relax}%
  \Gscale@div\@tempb{\linewidth}{\wd\pandoc@box}%
  \ifdim\@tempb\p@<\@tempa\p@\let\@tempa\@tempb\fi% select the smaller of both
  \ifdim\@tempa\p@<\p@\scalebox{\@tempa}{\usebox\pandoc@box}%
  \else\usebox{\pandoc@box}%
  \fi%
}
% Set default figure placement to htbp
\def\fps@figure{htbp}
\makeatother


% definitions for citeproc citations
\NewDocumentCommand\citeproctext{}{}
\NewDocumentCommand\citeproc{mm}{%
  \begingroup\def\citeproctext{#2}\cite{#1}\endgroup}
\makeatletter
 % allow citations to break across lines
 \let\@cite@ofmt\@firstofone
 % avoid brackets around text for \cite:
 \def\@biblabel#1{}
 \def\@cite#1#2{{#1\if@tempswa , #2\fi}}
\makeatother
\newlength{\cslhangindent}
\setlength{\cslhangindent}{1.5em}
\newlength{\csllabelwidth}
\setlength{\csllabelwidth}{3em}
\newenvironment{CSLReferences}[2] % #1 hanging-indent, #2 entry-spacing
 {\begin{list}{}{%
  \setlength{\itemindent}{0pt}
  \setlength{\leftmargin}{0pt}
  \setlength{\parsep}{0pt}
  % turn on hanging indent if param 1 is 1
  \ifodd #1
   \setlength{\leftmargin}{\cslhangindent}
   \setlength{\itemindent}{-1\cslhangindent}
  \fi
  % set entry spacing
  \setlength{\itemsep}{#2\baselineskip}}}
 {\end{list}}
\usepackage{calc}
\newcommand{\CSLBlock}[1]{\hfill\break\parbox[t]{\linewidth}{\strut\ignorespaces#1\strut}}
\newcommand{\CSLLeftMargin}[1]{\parbox[t]{\csllabelwidth}{\strut#1\strut}}
\newcommand{\CSLRightInline}[1]{\parbox[t]{\linewidth - \csllabelwidth}{\strut#1\strut}}
\newcommand{\CSLIndent}[1]{\hspace{\cslhangindent}#1}

\ifLuaTeX
\usepackage[bidi=basic]{babel}
\else
\usepackage[bidi=default]{babel}
\fi
% get rid of language-specific shorthands (see #6817):
\let\LanguageShortHands\languageshorthands
\def\languageshorthands#1{}


\setlength{\emergencystretch}{3em} % prevent overfull lines

\providecommand{\tightlist}{%
  \setlength{\itemsep}{0pt}\setlength{\parskip}{0pt}}



 


\AtBeginDocument{\hypersetup{linkcolor=black}}
% ---------- Paquetes tipograficos y microajustes ----------
\usepackage{microtype}        % mejor interletrado y justificado
\usepackage{csquotes}         % comillas tipograficas
\usepackage{iftex}
\ifPDFTeX\else
  % Seleccion de fuentes solo cuando XeTeX/LuaTeX esta disponible.
  \IfFontExistsTF{TeX Gyre Termes}{
    \setmainfont{TeX Gyre Termes}
  }{
    \IfFontExistsTF{Latin Modern Roman}{\setmainfont{Latin Modern Roman}}{}
  }
  \IfFontExistsTF{TeX Gyre Heros}{
    \setsansfont{TeX Gyre Heros}
  }{
    \IfFontExistsTF{Latin Modern Sans}{\setsansfont{Latin Modern Sans}}{}
  }
  \IfFontExistsTF{TeX Gyre Cursor}{
    \setmonofont{TeX Gyre Cursor}
  }{
    \IfFontExistsTF{Latin Modern Mono}{\setmonofont{Latin Modern Mono}}{}
  }
\fi
\usepackage{enumitem}         % listas compactas
\setlist{itemsep=.2em, topsep=.2em}

% ---------- KOMA-Script: estilo de titulos y espaciado ----------
\KOMAoptions{
  headings=big,
  parskip=half,
  fontsize=12pt,
  appendixprefix=true
}
\setlength{\parindent}{1.5em}
\setlength{\parskip}{0.6em}
\usepackage{indentfirst}
\usepackage{setspace}         % Control del interlineado del documento
\setstretch{1.5}

% ---------- Encabezados y pies (scrlayer-scrpage) ----------
\usepackage[automark,headsepline]{scrlayer-scrpage}
\clearpairofpagestyles
\automark[chapter]{chapter}
\ihead{\pagemark}
\ohead{\itshape\headmark}
\setheadsepline{0.4pt}
\renewcommand*{\chaptermarkformat}{}
\renewcommand*{\chapterpagestyle}{scrheadings}
\pagestyle{scrheadings}
\makeatletter
\let\ps@plain\ps@scrheadings
\makeatother
\cfoot{}

% ---------- Hipervinculos mas sobrios ----------
\usepackage{hyperref}
\hypersetup{
  colorlinks=true,
  linkcolor=black,
  citecolor=[rgb]{0.05,0.2,0.5},
  urlcolor=[rgb]{0.05,0.2,0.5},
  pdfauthor={\@author},
  pdftitle={\@title}
}

% ---------- Leyendas de figuras/tablas ----------
\usepackage[labelfont=bf,textfont=it]{caption}
\captionsetup{
  skip=8pt
}

% ---------- Tabla de contenidos ----------
\KOMAoptions{toc=graduated}
\RedeclareSectionCommand[tocnumwidth=3em]{chapter}
\RedeclareSectionCommand[tocindent=3.25em,tocnumwidth=2.8em]{section}
\RedeclareSectionCommand[tocindent=6.5em,tocnumwidth=2.5em]{subsection}
\makeatletter
\newcommand*{\tocdotfill}{\leavevmode\leaders\hbox to .6em{\hss.\hss}\hfill}
\makeatother
\RedeclareSectionCommand[toclinefill=\tocdotfill]{chapter}
\RedeclareSectionCommand[toclinefill=\tocdotfill]{section}
\RedeclareSectionCommand[toclinefill=\tocdotfill]{subsection}
\renewcommand*{\contentsname}{Tabla de contenido}
\setkomafont{chapterentry}{\normalfont}
\setkomafont{chapterentrypagenumber}{\normalfont}

% ---------- Viudas/Huerfanas y cortes de pagina ----------
\clubpenalty=10000
\widowpenalty=10000
\displaywidowpenalty=10000

% ---------- Entorno abstract para scrbook ----------
\providecommand{\abstractname}{Resumen}
\makeatletter
\@ifundefined{abstract}{
  \newenvironment{abstract}{
    \cleardoublepage
    \thispagestyle{plain}
    \null\vfill
    \begin{center}
      {\bfseries\Large \abstractname\par}
    \end{center}\vspace{1em}
    \begingroup
  }{
    \par\endgroup
    \vfill\null
    \cleardoublepage
  }
}{}
\makeatother

% ---------- Soporte de subtitulo desde YAML ----------
\makeatletter
\providecommand{\subtitle}[1]{\gdef\@subtitle{#1}}
\providecommand{\@subtitle}{}
\makeatother

\makeatletter
\newcommand{\facso@iflist}[4]{%
  \begingroup
    \def\addvspace##1{}%
    \def\numberline##1{##1}%
    \def\facso@target{#2}%
    \newcount\facso@entrycount
    \facso@entrycount=0\relax
    \long\def\contentsline##1##2##3{%
      \def\facso@this{##1}%
      \ifx\facso@this\facso@target
        \advance\facso@entrycount\@ne
      \fi
    }%
    \InputIfFileExists{\jobname.#1}{}{}%
  \endgroup
  \ifnum\facso@entrycount>0\relax
    #3%
  \else
    #4%
  \fi}
\renewcommand{\listoffigures}{%
  \facso@iflist{lof}{figure}{%
    \cleardoublepage
    \phantomsection
    \addchap*{Listado de Figuras}%
    \thispagestyle{scrheadings}%
    \markboth{Listado de Figuras}{Listado de Figuras}%
    \addcontentsline{toc}{chapter}{Listado de Figuras}%
    \@starttoc{lof}%
    \cleardoublepage
    \automark[chapter]{chapter}%
  }{}}
\renewcommand{\listoftables}{%
  \facso@iflist{lot}{table}{%
    \cleardoublepage
    \phantomsection
    \addchap*{Listado de Tablas}%
    \thispagestyle{scrheadings}%
    \markboth{Listado de Tablas}{Listado de Tablas}%
    \addcontentsline{toc}{chapter}{Listado de Tablas}%
    \@starttoc{lot}%
    \cleardoublepage
    \automark[chapter]{chapter}%
  }{}}
\makeatother

% --- Desactivar portada y abstract automaticos de Pandoc (PDF) ---
\AtBeginDocument{\let\maketitle\relax}
\renewenvironment{abstract}{}{}
\providecommand{\appendixname}{}
\providecommand{\appendixtocname}{}
\providecommand{\appendixpagename}{}
\renewcommand*{\appendixname}{Anexo}
\renewcommand*{\appendixtocname}{Anexos}
\renewcommand*{\appendixpagename}{Anexos}
\usepackage{booktabs}
\usepackage{longtable}
\usepackage{array}
\usepackage{multirow}
\usepackage{wrapfig}
\usepackage{float}
\usepackage{colortbl}
\usepackage{pdflscape}
\usepackage{tabu}
\usepackage{threeparttable}
\usepackage{threeparttablex}
\usepackage[normalem]{ulem}
\usepackage{makecell}
\usepackage{xcolor}
\makeatletter
\@ifpackageloaded{tcolorbox}{}{\usepackage[skins,breakable]{tcolorbox}}
\@ifpackageloaded{fontawesome5}{}{\usepackage{fontawesome5}}
\definecolor{quarto-callout-color}{HTML}{909090}
\definecolor{quarto-callout-note-color}{HTML}{0758E5}
\definecolor{quarto-callout-important-color}{HTML}{CC1914}
\definecolor{quarto-callout-warning-color}{HTML}{EB9113}
\definecolor{quarto-callout-tip-color}{HTML}{00A047}
\definecolor{quarto-callout-caution-color}{HTML}{FC5300}
\definecolor{quarto-callout-color-frame}{HTML}{acacac}
\definecolor{quarto-callout-note-color-frame}{HTML}{4582ec}
\definecolor{quarto-callout-important-color-frame}{HTML}{d9534f}
\definecolor{quarto-callout-warning-color-frame}{HTML}{f0ad4e}
\definecolor{quarto-callout-tip-color-frame}{HTML}{02b875}
\definecolor{quarto-callout-caution-color-frame}{HTML}{fd7e14}
\makeatother
\makeatletter
\@ifpackageloaded{bookmark}{}{\usepackage{bookmark}}
\makeatother
\makeatletter
\@ifpackageloaded{caption}{}{\usepackage{caption}}
\AtBeginDocument{%
\ifdefined\contentsname
  \renewcommand*\contentsname{Tabla de contenidos}
\else
  \newcommand\contentsname{Tabla de contenidos}
\fi
\ifdefined\listfigurename
  \renewcommand*\listfigurename{Listado de Figuras}
\else
  \newcommand\listfigurename{Listado de Figuras}
\fi
\ifdefined\listtablename
  \renewcommand*\listtablename{Listado de Tablas}
\else
  \newcommand\listtablename{Listado de Tablas}
\fi
\ifdefined\figurename
  \renewcommand*\figurename{Figura}
\else
  \newcommand\figurename{Figura}
\fi
\ifdefined\tablename
  \renewcommand*\tablename{Tabla}
\else
  \newcommand\tablename{Tabla}
\fi
}
\@ifpackageloaded{float}{}{\usepackage{float}}
\floatstyle{ruled}
\@ifundefined{c@chapter}{\newfloat{codelisting}{h}{lop}}{\newfloat{codelisting}{h}{lop}[chapter]}
\floatname{codelisting}{Listado}
\newcommand*\listoflistings{\listof{codelisting}{Listado de Listados}}
\makeatother
\makeatletter
\makeatother
\makeatletter
\@ifpackageloaded{caption}{}{\usepackage{caption}}
\@ifpackageloaded{subcaption}{}{\usepackage{subcaption}}
\makeatother
\usepackage{bookmark}
\IfFileExists{xurl.sty}{\usepackage{xurl}}{} % add URL line breaks if available
\urlstyle{same}
\hypersetup{
  pdftitle={Estudio de Notas de la Facultad de Ciencias Sociales},
  pdfauthor={Tomás Urzúa Moreno; Juan Carlos Castillo; Daniel Miranda},
  pdflang={es},
  colorlinks=true,
  linkcolor={Maroon},
  filecolor={Maroon},
  citecolor={Blue},
  urlcolor={Blue},
  pdfcreator={LaTeX via pandoc}}


\title{Estudio de Notas de la Facultad de Ciencias Sociales}
\usepackage{etoolbox}
\makeatletter
\providecommand{\subtitle}[1]{% add subtitle to \maketitle
  \apptocmd{\@title}{\par {\large #1 \par}}{}{}
}
\makeatother
\subtitle{Diagnóstico y análisis de las notas de estudiantes de
pregrado}
\author{Tomás Urzúa Moreno \and Juan Carlos Castillo \and Daniel
Miranda}
\date{18 de junio de 2026}
\begin{document}
\frontmatter
\maketitle

% --- title-pdf.tex personalizado para portada formal ---
\makeatletter
\providecommand{\subtitle}[1]{\gdef\@subtitle{#1}}
\providecommand{\@subtitle}{}
\providecommand{\frontmattercontext}{}
\providecommand{\advisorname}{}
\providecommand{\advisorlabel}{FACSO - Universidad de Chile}
\providecommand{\frontmatterlocation}{}
\InputIfFileExists{includes/cover-config.tex}{}{}

\newcommand{\PrintTitle}{%
  {\sffamily\bfseries\fontsize{24pt}{28pt}\selectfont \@title\par}%
}
\newcommand{\PrintSubtitle}{%
  \begingroup
  \edef\temp{\detokenize{\@subtitle}}%
  \ifx\temp\empty\relax
    % sin subtitulo
  \else
    {\sffamily\bfseries\large \@subtitle\par}%
  \fi
  \endgroup
}
\newcommand{\PrintAuthor}{%
  % Autores ocultos en la portada PDF
}
\newcommand{\PrintDate}{%
  {\small \@date\par}%
}
\makeatother

\begin{titlepage}
\thispagestyle{empty}
\begin{center}
\vspace*{10mm}

% Logo o imagen institucional

\includegraphics[width=0.35\textwidth]{assets/cover.jpg}\par
\vspace{12mm}

% Titulo y subtitulo
\PrintTitle
\vspace{6mm}
\PrintSubtitle

\vspace{15mm}
\begingroup
\edef\temp{\detokenize{\frontmattercontext}}%
\ifx\temp\empty\relax
  % sin contexto adicional
\else
  {\normalsize \frontmattercontext\par}
\fi
\endgroup

\vspace{18mm}
\PrintAuthor

\vspace{16mm}
\begingroup
\edef\temp{\detokenize{\advisorname}}%
\ifx\temp\empty\relax
  % sin tutor
\else
  \rule{0.45\textwidth}{0.4pt}\par
  {\small \advisorlabel\ \advisorname\par}
\fi
\endgroup

\vfill
\begingroup
\edef\temp{\detokenize{\frontmatterlocation}}%
\ifx\temp\empty\relax
  % sin ubicacion
\else
  {\small \frontmatterlocation\par}
\fi
\endgroup
\PrintDate

\end{center}
\end{titlepage}

% Preliminares (numeros romanos) y estilo simple
\frontmatter
\pagestyle{scrheadings}

\cleardoublepage
\thispagestyle{empty}
\begin{center}
\setlength{\fboxsep}{10pt}
\fcolorbox{gray!60}{gray!10}{%
\begin{minipage}{0.92\textwidth}
\textbf{Esta investigación fue desarrollada por Tomás Urzúa, sociólogo y estudiante del Magister en Ciencias Sociales del Departamento de Sociología de la Universidad de Chile.} Contó con la asesoría de los profesores Juan Carlos Castillo (Departamento de Sociología), y Daniel Miranda (Departamento de Psicología), de la Universidad de Chile.

El informe se encuentra disponible en versión web en https://data-soc.github.io/estudio-notas/ y en formato PDF en https://data-soc.github.io/estudio-notas/estudio-notas.pdf.

\textbf{Citar de la siguiente manera:}

\begin{itemize}
\tightlist
\item \emph{Urzúa, T., Castillo, J-C \& Miranda, D. (2026). Estudio de Notas: Análisis de las calificaciones de los estudiantes de la Facultad de Ciencias Sociales. Informe técnico. Dirección de Investigación, Facultad de Ciencias Sociales, Universidad de Chile. Disponible en https://data-soc.github.io/estudio-notas/estudio-notas.pdf}
\end{itemize}
\end{minipage}}
\end{center}
\cleardoublepage

\setcounter{tocdepth}{2} % (o 1/3 segun prefieras)
\tableofcontents
\listoffigures
\listoftables


\mainmatter
\bookmarksetup{startatroot}

\chapter*{Resumen ejecutivo}\label{resumen-ejecutivo}
\addcontentsline{toc}{chapter}{Resumen ejecutivo}

\markboth{Resumen ejecutivo}{Resumen ejecutivo}

\small\textbf{Palabras clave:} Notas, Facultad de Ciencias Sociales,
Factores asociados \normalsize

A pesar de que una de las funciones principales de una institución
educativa como es la Universidad de Chile es formar profesionales
competentes, la forma en que se evalúan estas competencias y los
aprendizajes rara vez ha sido analizada en profundidad. Sabemos poco
sobre si las notas que se otorgan a los estudiantes reflejan realmente
su rendimiento académico y si permiten distinguir entre estudiantes que
sobresalen académicamente de aquellos que necesitan apoyo para cumplir
con los aprendizajes esperados. En este contexto, el presente informe se
propone analizar las notas de los estudiantes de la Facultad de Ciencias
Sociales (FACSO) con el objetivo de describir la distribución de las
calificaciones, identificar patrones y tendencias a lo largo del tiempo,
y explorar los posibles factores asociados al rendimiento académico.
Este análisis permitirá comprender mejor el sistema de evaluación actual
y proponer recomendaciones para fortalecer la calidad educativa en la
facultad.

La literatura nacional e internacional sobre calificaciones en distintos
sistemas educativos ha señalado un fenómeno conocido como
\textbf{``inflación de notas''}, que se refiere a la tendencia de las
calificaciones a aumentar con el tiempo sin un correspondiente aumento
en el rendimiento académico real de los estudiantes. Este fenómeno ha
sido documentado en diversos contextos educativos y se ha asociado a
factores como cambios en los criterios de evaluación, presiones
institucionales para mejorar los indicadores de rendimiento, y la
influencia de factores socioeconómicos y demográficos en el desempeño
académico. Frente a este escenario, el presente estudio se propone
analizar si existe evidencia de inflación de notas en la Facultad de
Ciencias Sociales, y en caso afirmativo, identificar los factores que
podrían estar contribuyendo a este fenómeno. Para ello, se analizarán
las calificaciones de los estudiantes a lo largo del tiempo, se
compararán las tendencias entre diferentes carreras y se explorarán las
posibles asociaciones con variables como el género, el origen escolar,
el nivel socioeconómico, el tipo de curso y el tamaño del curso.

La interrogante inicial es aparentemente sencilla: \textbf{¿cuál es el
promedio de notas de la facultad?} Si bien todos estamos familiarizados
con la idea del promedio de notas, esta pregunta no es tan simple de
responder. Para calcular un promedio de notas representativo, es
necesario considerar una serie de factores que pueden influir en la
distribución de las calificaciones, como el número de estudiantes, el
número de cursos, la variabilidad entre carreras, y los cambios a lo
largo del tiempo. Además, es importante analizar no solo el promedio
general, sino también cómo se distribuyen las calificaciones dentro de
cada carrera y cómo estas han evolucionado en los últimos años. Por lo
tanto, el presente estudio se propone abordar esta pregunta desde una
perspectiva integral, analizando no solo el promedio de notas a nivel
facultad, sino también explorando las tendencias y patrones que emergen
al desagregar los datos por carrera y por año de ingreso.

El análisis se estructura en dos secciones principales, según distintas
unidades de análisis: \textbf{cursos} y \textbf{estudiantes}. Partiendo
por la sección de estudiantes, se analizarán las calificaciones a nivel
individual, explorando cómo estas se relacionan con variables
sociodemográficas como el género, el origen escolar y el nivel
socioeconómico. En la sección de notas por cursos, se analizarán las
calificaciones a nivel de cada curso impartido en la facultad,
identificando tendencias a lo largo del tiempo y comparando entre cursos
obligatorios y electivos. A lo largo del informe, se presentarán los
resultados más relevantes de cada sección, seguidos de una discusión que
integrará los hallazgos y propondrá recomendaciones para mejorar el
sistema de evaluación en la Facultad de Ciencias Sociales.

Los datos analizados provienen de las bases de datos de UCampus, que
incluyen información sobre las calificaciones de los estudiantes, así
como variables sociodemográficas y académicas relevantes. El análisis se
ha realizado utilizando técnicas estadísticas descriptivas y
correlacionales, con el objetivo de identificar patrones y tendencias en
la distribución de las calificaciones. Además, se han aplicado métodos
de visualización de datos para facilitar la interpretación de los
resultados y comunicar de manera efectiva los hallazgos a la comunidad
académica. Junto con los gráficos provistos, se ha generado un
\textbf{visualizador de notas por cursos}, que es una plataforma
interactiva que permite explorar los datos de manera dinámica,
facilitando la comprensión de las tendencias y patrones identificados en
el análisis.

\textbf{Principales resultados sección estudiantes (N = 1963,
2020-2024)}

\begin{itemize}
\item
  Los análisis iniciales en base a datos de UCampus (con el promedio
  general de cada estudiante) indican que el promedio de notas general a
  nivel facultad es \textbf{6.4} entre 2020 y 2024, con una desviación
  estándar de 0.33. Esto quiere decir que la mayoría de los estudiantes
  obtienen calificaciones cercanas a este promedio, con una variabilidad
  relativamente baja.
\item
  Al considerar una segunda base de datos con información de todos los
  cursos, incluyendo más estudiantes y los cursos reprobados, la nota
  promedio baja a un \textbf{6.1}, con una desviación estándar de 0.6.
  Esto sugiere que al incluir los cursos reprobados, la distribución de
  las calificaciones se vuelve más amplia, reflejando una mayor
  variabilidad en el rendimiento académico de los estudiantes. Para los
  análisis se decide trabajar con esta segunda base de datos, ya que
  permite una representación más completa y precisa del rendimiento
  académico de los estudiantes, incluyendo tanto los éxitos como los
  desafíos que enfrentan en su trayectoria académica.
\item
  En relación con el género, se observa una tendencia consistente en
  todas las carreras: las mujeres obtienen, en promedio, mejores
  calificaciones que los hombres. Esta diferencia se mantiene
  independientemente del año de ingreso y alcanza su mayor magnitud en
  Trabajo Social, con una brecha aproximada de 0.54 puntos a favor de
  las estudiantes.
\item
  Respecto al origen escolar, se identifican patrones diferenciados
  entre carreras. En Antropología, Educación Parvularia y Sociología se
  observa una estructura escalonada según el tipo de establecimiento
  escolar: los estudiantes provenientes de colegios municipales
  presentan las notas más bajas, seguidos por quienes egresaron de
  particulares subvencionados, mientras que los de particulares pagados
  obtienen los promedios más altos. En contraste, en Psicología y
  Trabajo Social los estudiantes de colegios particulares subvencionados
  presentan el mejor rendimiento promedio, superando incluso a los
  provenientes de colegios particulares pagados. Asimismo, en las
  cohortes más recientes se observa una creciente brecha entre
  estudiantes provenientes de colegios municipales y aquellos de
  establecimientos particulares.
\item
  Al analizar el nivel socioeconómico, en Antropología, Psicología y
  Sociología se aprecia una relación positiva entre estatus
  socioeconómico y rendimiento académico. Sin embargo, esta tendencia no
  se replica en Educación Parvularia ni en Trabajo Social, donde se
  observan patrones distintos: en la primera, los estudiantes de nivel
  medio obtienen las notas más altas, mientras que en la segunda los
  estudiantes de nivel socioeconómico alto presentan los promedios más
  bajos.
\end{itemize}

\textbf{Principales resultados sección cursos (N = 2291, 2018-2025)}

\begin{itemize}
\item
  El análisis de las calificaciones a nivel de cursos revela que el
  promedio de notas a nivel facultad ha experimentado un aumento
  significativo durante el período analizado, pasando desde medianas
  cercanas a 5,5--5,7 hacia valores sobre 6,1, lo que sugiere un
  incremento generalizado en las calificaciones. Este hallazgo es
  consistente con la hipótesis de inflación de notas, sugiriendo que los
  estándares de evaluación podrían haberse flexibilizado durante este
  período, especialmente durante la pandemia.
\item
  Una variable relevante a nivel de curso es la \textbf{tasa de
  reprobación}, que en promedio es de un \textbf{3\%} a nivel facultad,
  lo que indica que la mayoría de los estudiantes aprueban sus cursos.
  En relación con esto, la dispersión de las calificaciones a nivel
  facultad ha seguido una tendencia decreciente en el período analizado,
  evidenciando una progresiva homogeneización de los resultados
  evaluativos.
\item
  Psicología y Educación Parvularia son las carreras más estables y con
  tasas sostenidamente altas, rozando el 99\% en varios años.
  Antropología, en cambio, muestra la trayectoria más volátil y una
  tendencia levemente decreciente, partiendo cerca del 100\% en 2018 y
  situándose en torno al 96\% hacia 2024--2025. Por su parte, Sociología
  es la carrera con mayor variabilidad entre años y registra los valores
  más bajos del conjunto, con una caída en 2019 (\textasciitilde94\%) y
  fluctuaciones persistentes que la ubican por debajo del resto de
  carreras. Trabajo Social presenta un comportamiento intermedio, con
  ciertas variaciones pero relativamente estable.
\item
  Los cursos electivos registran calificaciones consistentemente más
  elevadas que los obligatorios, lo que sugiere que la motivación o
  afinidad asociada a la elección voluntaria de estos cursos podría
  operar como un factor que contribuye a elevar los promedios
  individuales.
\end{itemize}

\textbf{Conclusiones}

El informe presenta evidencia a favor de la hipótesis de inflación de
notas al notar el alza de las calificaciones durante la pandemia, nivel
que aún se mantiene por sobre los promedios de calificaciones
prepandémicos. Esto da cuenta de un fenómeno global en la educación
superior en el que los estándares de evaluación han disminuido su
exigencia durante los últimos años.

Consistente con la inflación de notas, las calificaciones otorgadas
hacen difícil distinguir entre estudiantes que sobresalen académicamente
de aquellos que necesitan apoyo para cumplir con los aprendizajes
esperados debido a la baja variabilidad de las notas que obtienen los
estudiantes. Es más, las calificaciones han tendido a homogeneizarse con
el paso del tiempo, dificultando aún más la tarea de diferenciar grupos
de estudiantes a partir de su rendimiento académico.

A lo largo del tiempo, el promedio de aprobación de cursos a nivel
facultad es de un 97\%, demostrando que reprobar en FACSO es una
anomalía. Con ello, se da cuenta de la baja exigencia que implica
adquirir las credenciales de aprendizajes esperados en los ramos que
imparte la facultad. Factores socioeconómicos (nivel socioeconómico y
tipo de colegio), sociodemográficos (género) e institucionales (tipo de
curso y tamaño del curso) permiten en parte explicar la distribución de
las notas a nivel facultad y a nivel de cada carrera.

En el marco de la excelencia académica que identifica a la Universidad
de Chile, se enfatiza abrir la discusión sobre los procedimientos
evaluativos actuales que se llevan en cada carrera, con el objetivo de
lograr discernir entre aquellos estudiantes que tienen rendimientos por
sobre el promedio y los que necesitan políticas de acompañamiento para
lograr una formación académico-profesional exitosa. Se proponen tres
investigaciones para poder avanzar en este tema: complementar el Estudio
de Notas con información de profesores, tales como el dictado de cursos
y la evaluación docente (1), aplicar un estudio comparativo de los
criterios de evaluación de los ramos impartidos en cada carrera (2) y
replicar el Estudio de Notas en otras facultades de la Universidad de
Chile (3).

Este informe se ha elaborado siguiendo principios de ciencia abierta,
aplicando prácticas que promueven la transparencia y la reproducibilidad
de la investigación. Todos los datos, así como el procesamiento de los
mismos y los códigos generados para los análisis, se encuentran
disponibles en el
\href{https://github.com/tomasurzuam/estudio-notas}{repositorio GitHub
asociado al estudio}.

\bookmarksetup{startatroot}

\chapter{Introducción}\label{introducciuxf3n}

La excelencia académica ha sido un principio constitutivo de la
tradición intelectual y la identidad cultural de la Universidad de
Chile, orientando su modelo educativo hacia la formación de
profesionales con vocación de aporte público y desarrollo pleno de sus
talentos. Este principio se expresa en la promoción de aprendizajes
basados en competencias, entendidas como un conjunto articulado de
saberes que permiten a los y las egresadas tomar decisiones
fundamentadas en diversos ámbitos de desempeño
(\citeproc{ref-universidaddechile_modelo_2024}{Universidad de Chile,
2024}). En coherencia con su compromiso público, la Universidad reconoce
la heterogeneidad de las trayectorias de origen de su estudiantado y
promueve la igualdad de oportunidades mediante políticas de inclusión y
acompañamiento académico, las cuales apuntan a favorecer la integración
a la vida universitaria, la titulación oportuna y la inserción laboral.
De este modo, la excelencia académica y la inclusión se configuran como
una condición esencial del derecho a la educación y como un pilar para
la construcción de una sociedad democrática
(\citeproc{ref-universidaddechile_modelo_2024}{Universidad de Chile,
2024}).

En la agenda investigativa de la educación superior, una de las
problemáticas que ha cobrado creciente relevancia es la asignación de
notas, en tanto el rendimiento académico ha sido históricamente uno de
los elementos constitutivos del sistema universitario y el principal
indicador para evaluar el logro de los aprendizajes de los estudiantes
(\citeproc{ref-lynch_learning_2017}{Lynch \& Hennessy, 2017}). En este
sentido, las calificaciones operan como un mecanismo de
credencialización, al evidenciar si los contenidos impartidos en una
carrera universitaria han sido efectivamente aprehendidos. Durante los
últimos años, sin embargo, se ha develado un alza sostenida de las
calificaciones académicas, patrón que resulta consistente en
instituciones universitarias a nivel internacional
(\citeproc{ref-achen_what_2009}{Achen \& Courant, 2009};
\citeproc{ref-jewell_whose_2011}{Jewell et~al., 2011}). Este fenómeno ha
generado un debate en torno a las razones que explican dicho incremento,
ilustrado en la disyuntiva sobre si los estudiantes se han vuelto más
competentes o si, por el contrario, las notas han experimentado un
proceso de inflación (\citeproc{ref-jephcote_grade_2021}{Jephcote
et~al., 2021}). En tal contexto, una manera de operacionalizar la
inflación de notas ha sido mediante la calificación de los estudiantes.
Frente a estas problemáticas, la literatura no ha alcanzado un consenso
claro respecto de sus causas: mientras algunas investigaciones sostienen
que, con el paso del tiempo, ciertos departamentos ---particularmente en
ciencias sociales y humanidades--- han reducido la exigencia de sus
evaluaciones, contribuyendo así al aumento de los promedios de notas
(\citeproc{ref-hernandez-julian_measuring_2016}{Hernández-Julián \&
Looney, 2016}; \citeproc{ref-lipnevich_what_2020}{Lipnevich et~al.,
2020}), otros hallazgos atribuyen este fenómeno a un incremento efectivo
en la calidad académica de los estudiantes
(\citeproc{ref-hernandez-julian_measuring_2016}{Hernández-Julián \&
Looney, 2016}).

La elevación sistemática de las calificaciones tiene consecuencias
relevantes a nivel institucional y educativo. En el plano institucional,
dificulta identificar a estudiantes con mayores necesidades de apoyo,
obstaculizando la focalización de políticas de acompañamiento académico,
y también deteriora la confiabilidad de las notas como indicador de
competencias al proyectar una imagen sobrevalorada del desempeño
(\citeproc{ref-schorr_addressing_2025}{Schorr, 2025}). En el ámbito
educativo, obtener buenas calificaciones independientemente de las
exigencias del curso desincentiva la asistencia a clases, genera
percepciones de injusticia entre estudiantes con distintos niveles de
compromiso y aumenta el estrés académico, ya que cualquier resultado
bajo el estándar inflado es percibido como un fracaso
(\citeproc{ref-schorr_addressing_2025}{Schorr, 2025}).

Los factores asociados a las calificaciones universitarias pueden
encontrar diversas razones. Jewell et~al.
(\citeproc{ref-jewell_whose_2011}{2011}) evidenció que algunos
profesores tienden a asignar mejores notas con tal de recibir
evaluaciones docentes positivas, lo que puede tener retribuciones
directas en procesos de promoción y/o permanencia. Por otro lado, desde
una racionalidad económica, calificar con bajas notas requiere mayor
inversión de tiempo y energía, pues se debe entregar una
retroalimentación más detallada en comparación con un trabajo mejor
elaborado. Entonces, los docentes podrían estar evaluando con menor
exigencia a los estudiantes con el fin de reducir su tiempo en tareas de
revisión y, así, maximizarlo para otras labores. Una última posible
causa se puede encontrar en la cultura de las adecuaciones que ha
caracterizado a los estudiantes, sobre todo el período posterior a la
pandemia. Las adecuaciones curriculares, que son casi en su totalidad
por salud mental, entregan una amplia flexibilidad para cumplir con las
responsabilidades académicas, llegando al punto de que los
universitarios pueden ausentarse a las evaluaciones regulares de los
cursos y solamente presentarse al examen final, o bien que la asistencia
obligatoria pierda este carácter. Si bien estas medidas pretenden
brindar una estabilidad psicoemocional a los estudiantes, también
podrían estar afectando su rendimiento sin que esto se vea reflejado en
sus calificaciones.

En el segundo semestre del 2025, la Facultad de Ciencias Sociales
(FACSO) de la Universidad de Chile se planificó un estudio de notas que
contemplaba, inicialmente, realizar un diagnóstico general en cuanto al
rendimiento académico del cuerpo estudiantil, con el objetivo de
identificar y apoyar a aquellos estudiantes que presentan mayores
dificultades en el logro de resultados. Este
\href{https://tomasurzuam.github.io/estudio-notas/output/reports/ReporteA_2025.html}{análisis
preliminar} se llevó a cabo por la Dirección de Pregrado, realizado a
partir de los promedios finales de estudiantes regulares entre 2020 y
2025\footnote{Este promedio es una nota única que registra UCampus
  considerando todos los períodos académicos que el estudiante terminó.
  Así, para un estudiante que ingresó el 2025 se refleja sólo el primer
  semestre, mientras que para un estudiante que ingresó el 2021 su nota
  considera los promedios de 9 semestres.}, considerando las siete
carreras que imparte FACSO (N = 2578). Los principales hallazgos fueron,
en primer lugar, que el promedio de notas de FACSO sería \textbf{6,4}.
Si bien este alto promedio podría interpretarse como que la mayoría de
los estudiantes tiene un rendimiento sobresaliente, también es dable
pensar que estamos ante un escenario de inflación de notas, lo que tiene
una serie de implicancias. En primer lugar, no permite discriminar entre
aquellos estudiantes que tienen un rendimiento académico destacado y
aquellos que podrían estar enfrentando dificultades, lo que dificulta la
implementación de estrategias de apoyo efectivas. Segundo, los
estudiantes podrían estar recibiendo una retroalimentación inadecuada
sobre su desempeño, lo que podría afectar su motivación y su desarrollo
académico a largo plazo. Tercero, la inflación de notas puede generar
una percepción de injusticia entre los estudiantes, ya que aquellos que
se esfuerzan por obtener buenas calificaciones podrían sentirse
desmotivados al ver que otros reciben calificaciones similares sin el
mismo nivel de esfuerzo. En cuarto lugar, en este escenario la
asistencia a clases también se vuelve irrelevante, ya que las
calificaciones no dependen de asistir o no. Por último, este fenómeno
podría estar contribuyendo a un aumento del estrés académico, ya que los
estudiantes podrían sentir la presión de mantener un alto rendimiento en
un entorno donde las expectativas son cada vez más elevadas.

Debido a ello, se le solicitó a la Dirección de Investigación de FACSO
realizar un estudio detallado sobre las calificaciones en la facultad,
con el fin de abordar las interrogantes relacionadas a la distribución
de las calificaciones y los modos de evaluación. Considerando lo
anterior, el presente reporte plantea cuatro preguntas: ¿Existe la
inflación de notas en la Facultad de Ciencias Sociales? (1) ¿Las notas
permiten discriminar grupos de estudiantes en función del rendimiento?
(2) ¿Qué tan alto es el nivel de exigencia en cuánto a la aprobación de
asignaturas? (3) ¿Qué factores asociados dan cuenta de las notas, el
grado de discriminación y la tasa de aprobación? (4).

Para responder estas incógnitas, el estudio plantea dos objetivos
principales: (1) \textbf{describir y analizar el rendimiento académico
en las y los estudiantes de pregrado de FACSO sus cambios en el tiempo y
sus factores asociados}, y (2) \textbf{describir y analizar el
rendimiento académico de los cursos de pregrado de FACSO entre 2018 y
2025, sus variaciones en el tiempo y características asociadas} (2). Con
el fin de cumplir estos objetivos, este estudio propone una
operacionalización de la inflación de notas que va más allá del mero
promedio, integrando la dispersión de notas y el porcentaje de
aprobación.

El informe se estructurará a partir de las siguientes secciones:
Metodología, Resultados y Conclusiones. En la metodología se presenta la
información en detalle de los datos y técnicas estadísticas que se han
utilizado para desarrollar la investigación. Posteriormente, el apartado
de resultados se divide en un subapartado para estudiantes y otro para
cursos, los que contienen la totalidad de los análisis, estudiando la
situación de la Facultad, así como las diferencias por carrera y en el
tiempo. Por último, en la conclusión se resumen los principales
resultados de la investigación y se proponen lineamientos de acción para
profundizar en esta agenda.

\bookmarksetup{startatroot}

\chapter{Metodología}\label{metodologuxeda}

Este apartado se divide en dos secciones principales: Métodos para la
dimensión de estudiantes (1) y para la dimensión de cursos (2). Cada una
de ellas se compone de tres subsecciones. Primero, se presentan los
datos con los cuales se realiza el informe, detallando las bases de
datos disponibles desde un inicio, las tomas de decisiones para el
procesamiento de las mismas y la nueva generación de información. En
segundo lugar, se presentan las variables disponibles del estudio, junto
con su descripción y algunos estadísticos de las variables de interés.
Por último, se define la estrategia analítica a utilizar para elaborar
los resultados. Cabe mencionar que todos los códigos asociados al
procesamiento y análisis de los datos se encuentran disponibles en el
\href{https://github.com/tomasurzuam/estudio-notas}{repositorio GitHub
asociado al estudio}

\section{Sección estudiantes}\label{secciuxf3n-estudiantes}

\subsection{Datos}\label{datos}

Al comienzo de este estudio, se entregaron dos bases de datos
previamente construidas a partir de información de UCampus: la base
\texttt{Original\ general} y \texttt{Notas\ por\ curso}. La primera base
contiene información sobre estudiantes de todas las carreras de FACSO
entre los años 2020 y 2025. Cada estudiante tiene una nota que
representa el promedio de las notas de todos sus cursos. Además,
presenta variables de caracterización sociodemográficas y académicas,
las cuales serán detalladas más adelante. La segunda base tiene las
notas desagregadas por curso entre 2020 y 2024, es decir, posee las
calificaciones de todos los ramos que cursaron durante ese período. No
obstante, no cuenta con variables de caracterización ni un promedio
agregado como la base \texttt{Original\ general}.

Todas las variables de caracterización son propias de UCampus a
excepción de nivel socioeconómico, ya que esta variable fue obtenida
mediante una encuesta institucional de caracterización aplicada de
manera externa a la plataforma UCampus, y posteriormente fue integrada a
sus bases de datos. Esta variable refleja el nivel de vulnerabilidad de
la comuna de origen del estudiante, con valores que oscilan entre 0 y 8,
donde los valores más bajos indican mayor vulnerabilidad socioeconómica.
Con el propósito de facilitar su interpretación y análisis, la variable
fue recategorizada en tres niveles: bajo (0, 1 y 2), medio (3, 4 y 5) y
alto (6, 7 y 8). Las demás variables de caracterización tuvieron
procesamientos menores, tales como etiquetado y transformación a factor.

A partir de las dos bases que fueron dispuestas para trabajar, se generó
una \texttt{Base\ madre}, que es un dataframe creado a partir del merge
de las bases \texttt{Original\ general} y \texttt{Notas\ por\ curso},
uniéndose mediante el RUT del estudiante. Se decidió esto para tener una
base que contuviera la mayor cantidad de información posible para el
análisis, obteniendo las notas por curso y variables de caracterización
en una sola base de datos. Debido a la incompatibilidad de los años
presentes en ambas bases, se logró mantener información entre 2020 y
2024, perdiéndose el año 2025.

La base \texttt{Original\ general} contiene 2578 casos válidos. Por su
parte, \texttt{Notas\ por\ curso} tiene 3638 casos. Posterior al merge,
\texttt{Base\ madre} queda con un N total de 1981 individuos. La
disminución de casos en la base final se debe a la falta de información
disponible en ambas bases, es decir, hay estudiantes que están presentes
en la base \texttt{Original\ general} pero no en
\texttt{Notas\ por\ curso}, y viceversa.

Posterior a la generación de la \texttt{Base\ madre}, se crearon
múltiples bases filtrando por carrera, por lo que la estructura de los
datos es idéntica, lo cual aporta a la eficiencia de los análisis y
reproducibilidad del código. En este marco, todas las bases que parten
con \texttt{Notas} seguido del nombre de la carrera, fueron construidas
de esta forma. Todas las bases mencionadas tienen como unidad de
análisis a los estudiantes.

Para llevar a cabo los análisis, se utilizan las notas de la base
\texttt{Notas\ por\ curso}, construyéndose una nueva variable de
promedio de notas por estudiante, la cual se considera más fiable al
haber sido procesada por el mismo equipo de investigación. Para la
construcción de esta nueva variable de notas, se tomaron las siguientes
decisiones: (1) se recodificaron como NA los valores iguales a 0 y 1
debido a que este era su significado real. (2) Para calcular el
promedio, se consideraron la totalidad los ramos que cursó un
estudiante. En este contexto, si un estudiante cursó tres veces un ramo,
en el cual las primeras dos veces lo reprobó y la última lo aprobó, esas
tres notas se contemplan para el promedio, al igual que el resto de los
cursos finalizados por la persona. Por último, para determinar factores
que afectan las notas de los estudiantes, se recurre a variables de
caracterización de origen en la base \texttt{Original\ general}.

A partir de todas las decisiones de procesamiento tomadas, la muestra
analítica se constituyó en un total de 1981 cursos.

Todas las bases de datos se pueden encontrar en
\href{https://github.com/tomasurzuam/estudio-notas/tree/main/input/data}{este
repositorio}.

\begin{table}

\caption{\label{tbl-bbdd}Bases de datos disponibles}

\centering{

\centering
\resizebox{\ifdim\width>\linewidth\linewidth\else\width\fi}{!}{
\begin{tabular}{l|l|l|l|l}
\hline
Base de datos & Código de la base & N & Formato & Base creada\\
\hline
Original general & original.rds & 2578 & Wide & No\\
\hline
Notas por curso & notas\_cursos.rds & 291675 & Long & No\\
\hline
Notas por curso procesada & notas\_cursos\_proc.rds & 3638 & Wide & Sí\\
\hline
Base madre & base\_madre.rds & 1981 & Wide & Sí\\
\hline
Notas sociología & df\_socio.rds & 459 & Wide & Sí\\
\hline
Notas psicología & df\_psico.rds & 549 & Wide & Sí\\
\hline
Notas antropología & df\_antropo.rds & 468 & Wide & Sí\\
\hline
Notas trabajo social & df\_trabajo.rds & 294 & Wide & Sí\\
\hline
Notas educación parvularia & df\_parvularia.rds & 211 & Wide & Sí\\
\hline
\end{tabular}}

}

\end{table}%

\subsection{Variables}\label{variables}

\begin{table}

\caption{\label{tbl-var}Variables disponibles}

\centering{

\centering\begingroup\fontsize{6}{8}\selectfont

\begin{tabular}{>{\raggedright\arraybackslash}p{2.2cm}|>{\raggedright\arraybackslash}p{1.7cm}|>{\raggedright\arraybackslash}p{1.6cm}|>{\raggedright\arraybackslash}p{6cm}|>{\raggedright\arraybackslash}p{1.9cm}|>{\raggedright\arraybackslash}p{1.2cm}}
\hline
Variable & Codigo & Tipo de variable & Categorias de respuesta & Valores de respuesta & Variable creada\\
\hline
\cellcolor{gray!10}{RUT} & \cellcolor{gray!10}{id} & \cellcolor{gray!10}{Numerica} & \cellcolor{gray!10}{-} & \cellcolor{gray!10}{-} & \cellcolor{gray!10}{No}\\
\hline
Tipo de ingreso a la universidad & ingreso & Categorica & "ESTUDIANTES CIEGOS" "LICENCIADOS Y TITULADOS"        
  "TRANSFERENCIAS EXTERNAS" "EQUIDAD DE GENERO"              
  "DEPORTISTA DESTACADO" "EST. MEDIOS EXTRANJERO"         
  "BEA (BECA EXCELENCIA ACADEMICA)" "SIPEE"                          
  "PACE" "PSU" "PUEBLOS INDIGENAS" "DISCAPACIDAD" & - & No\\
\hline
\cellcolor{gray!10}{Sexo} & \cellcolor{gray!10}{sexo} & \cellcolor{gray!10}{Dicotomica} & \cellcolor{gray!10}{"Hombre" "Mujer"} & \cellcolor{gray!10}{0 | 1} & \cellcolor{gray!10}{No}\\
\hline
Nacionalidad & nacionalidad & Categorica & "Chilena"                      "Otra nacionalidad;Chilena"   
 "Alemana;Venezolana"           "Venezolana"                  
 "Chilena;Argentina"            "Ecuatoriana;Chilena"         
 "Peruana;Chilena"              "Otra nacionalidad;Argentina" 
 "Otra nacionalidad;Colombiana" "Peruana"                     
 "Ecuatoriana"                  "Francesa;Chilena"            
 "Boliviana"                    "Brasileña"                   
 "Colombiana"                   "Suiza;Chilena"               
 "Costarricense;Chilena"        "Italiana;Chilena"            
 "Argentina"                    "Uruguaya"                    
 "Chilena;Uruguaya" & - & No\\
\hline
Origen etnico & origen\_etnico & Categorica & "NINGUNA"                  "MAPUCHE"                 
  "AYMARA"                   NA                        
  "DIAGUITA"                 "QUECHUA"                 
  "SIN INFORMACION"          "AFRODESCENDIENTE CHILENO"
  "ATACAMEÑO o LICKANANTAI"  "RAPA NUI"                
\cellcolor{gray!10}{ "COLLA"} & \cellcolor{gray!10}{-} & \cellcolor{gray!10}{No}\\
\hline
Situacion academica & sit\_academica & Categorica & "MATRICULADO"                "POSTERGACION   1º SEMESTRE"
 "POSTERGACION ANUAL"         "TRANSFERENCIA INTERNA"     
 "RENUNCIA  A LA CARRERA" & - & No\\
\hline
Carrera & carrera & Categorica & "PSICOLOGIA"                        "ANTROPOLOGIA ARQUEOLOGIA"         
 "SOCIOLOGIA"                        "TRABAJO SOCIAL"                   
\cellcolor{gray!10}{ "PEDAGOGIA EN EDUCACION PARVULARIA"} & \cellcolor{gray!10}{-} & \cellcolor{gray!10}{No}\\
\hline
Cohorte & cohorte & Categorica & 2020 2021 2022 2023 2024 2025 & - & No\\
\hline
Tipo de colegio de procedencia & colegio & Categorica & "S/INF"                    "MUNICIPAL"               
\cellcolor{gray!10}{ "PARTICULAR SUBVENCIONADO" "PARTICULAR PAGADO"} & \cellcolor{gray!10}{-} & \cellcolor{gray!10}{No}\\
\hline
Año de egreso de enseñanza media & egreso\_emedia & Categorica & 2003 2006 2007 2013 2015 2016 2017 2019 2018 2009 2014 2020 2011 2008
 2012 2021 2002 2022 2001 2010 2023 1986 2005 & - & No\\
\hline
\cellcolor{gray!10}{Promedio de enseñanza media} & \cellcolor{gray!10}{promedio\_emedia} & \cellcolor{gray!10}{Categorica} & \cellcolor{gray!10}{-} & \cellcolor{gray!10}{-} & \cellcolor{gray!10}{No}\\
\hline
Nivel socioeconómico & nse & Ordinal & 0 5 4 7 3 8 2 6 1 & 0 | 1 | 2 | 3 | 4 | 5 | 6 | 7 | 8 & No\\
\hline
\cellcolor{gray!10}{Percentil en PAES o PSU} & \cellcolor{gray!10}{pct\_paes\_psu} & \cellcolor{gray!10}{Numerica} & \cellcolor{gray!10}{-} & \cellcolor{gray!10}{-} & \cellcolor{gray!10}{Sí}\\
\hline
Decil en PAES o PSU & decil\_paes\_psu & Numerica & - & - & Sí\\
\hline
Region de origen & region\_origen & Categorica & "REGION METROPOLITANA"          
 "REGION DEL BIOBÍO"              "REGIÓN DE ÑUBLE"               
 "REGION DE LOS LAGOS"            "REGION DE LA ARAUCANIA"        
 "REGION DE VALPARAÍSO"           "REGION DEL LIB. B. OHIGGINS"   
 "SIN INFORMACIÓN"                "REGION DE TARAPACA"            
 "REGION DE COQUIMBO"             "REGHON DE COQUIMBO"            
 "REGION DEL MAULE"               "REGION DE MAGALLANES"          
 "REGION DE AYSÉN"                "REGION DE ARICA Y PARINACOTA"  
 "REGION DE METROPOLITANA"        "REGION DE LOS RIOS"            
 "REGION DE ANTOFAGASTA"          "REGION DE IQUIQUE Y PARINACOTA"
 "REGION DE ATACAMA"              "REGION DE O'HIGGINS"           
\cellcolor{gray!10}{ "REGION DEL ÑUBLE"} & \cellcolor{gray!10}{-} & \cellcolor{gray!10}{No}\\
\hline
Cursos & prefijo "curso\_" seguido del nombre del curso & Numerica & - & - & No\\
\hline
\cellcolor{gray!10}{Promedio de notas por defecto} & \cellcolor{gray!10}{promedio\_ucampus} & \cellcolor{gray!10}{Numerica} & \cellcolor{gray!10}{-} & \cellcolor{gray!10}{-} & \cellcolor{gray!10}{No}\\
\hline
Promedio de notas calculado & promedio\_calculado & Numerica & - & - & Sí\\
\hline
\end{tabular}
\endgroup{}

}

\end{table}%

En la Tabla~\ref{tbl-var} se pueden observar todas las variables
disponibles del estudio esquemáticamente, presentándose su naturaleza,
codificación y categorías de respuesta. Para efectos del análisis, se
realizaron procesamientos a algunas variables, como es el caso de nivel
socioeconómico (NSE). En principio, esta variable poseía 9 categorías,
que iban del 0 al 8, pero se recategorizó de tal forma que quedaron tres
categorías: 0, 1 y 2 = Bajo; 3, 4, 5 = Medio; 6, 7 y 8 = Alto. Por otro
lado, para las bases centradas en estudiantes, se consideraron los años
de cohorte desde el 2021 al 2024, pues son los años en común que estaban
disponibles para la \texttt{Base\ madre}, permitiendo observar las notas
de cursos junto con caracterización sociodemográfica.

\begin{table}

\caption{\label{tbl-desc-num}Descriptivos variables numéricas}

\centering{

\centering\begingroup\fontsize{13}{15}\selectfont

\resizebox{\ifdim\width>\linewidth\linewidth\else\width\fi}{!}{
\begin{tabular}{>{\raggedright\arraybackslash}p{10em}|c|c|c|c|c|c}
\hline
\multicolumn{1}{c|}{\cellcolor[HTML]{952626}{\textcolor{white}{\textbf{ }}}} & \multicolumn{2}{c|}{\cellcolor[HTML]{952626}{\textcolor{white}{\textbf{Datos}}}} & \multicolumn{3}{c|}{\cellcolor[HTML]{952626}{\textcolor{white}{\textbf{Estadísticos Descriptivos}}}} & \multicolumn{1}{c}{\cellcolor[HTML]{952626}{\textcolor{white}{\textbf{Distribución}}}} \\
\cline{2-3} \cline{4-6} \cline{7-7}
\cellcolor[HTML]{952626}{\textcolor{white}{\textbf{Variable}}} & \cellcolor[HTML]{952626}{\textcolor{white}{\textbf{\% Casos Válidos}}} & \cellcolor[HTML]{952626}{\textcolor{white}{\textbf{\% Casos Perdidos}}} & \cellcolor[HTML]{952626}{\textcolor{white}{\textbf{Media}}} & \cellcolor[HTML]{952626}{\textcolor{white}{\textbf{Desviación Estándar}}} & \cellcolor[HTML]{952626}{\textcolor{white}{\textbf{Rango}}} & \cellcolor[HTML]{952626}{\textcolor{white}{\textbf{Histograma}}}\\
\hline
\textbf{Promedio calculado} & 100\% & 0\% & 5.92 & 0.58 & (1.5-6.84615384615385) & ▁▁▁▂▇\\
\hline
\textbf{Percentil prueba admisión} & 95.31\% & 4.69\% & 50.09 & 28.91 & (0.116550116550117-99.8834498834499) & ▇▇▇▇▇\\
\hline
\end{tabular}}
\endgroup{}

}

\end{table}%

\begin{table}

\caption{\label{tbl-desc-cat}Descriptivos variables categóricas}

\centering{

\centering
\caption{\label{tab:tbl-desc-cat}Estadísticas Descriptivas de Variables Categóricas}
\centering
\resizebox{\ifdim\width>\linewidth\linewidth\else\width\fi}{!}{
\fontsize{13}{15}\selectfont
\begin{tabular}[t]{>{\raggedright\arraybackslash}p{10em}|c|c|c|l}
\hline
\multicolumn{1}{c|}{\cellcolor[HTML]{952626}{\textcolor{white}{\textbf{ }}}} & \multicolumn{2}{c|}{\cellcolor[HTML]{952626}{\textcolor{white}{\textbf{Muestra}}}} & \multicolumn{2}{c}{\cellcolor[HTML]{952626}{\textcolor{white}{\textbf{Categorías}}}} \\
\cline{2-3} \cline{4-5}
\cellcolor[HTML]{952626}{\textcolor{white}{\textbf{Variable}}} & \cellcolor[HTML]{952626}{\textcolor{white}{\textbf{\% Casos válidos}}} & \cellcolor[HTML]{952626}{\textcolor{white}{\textbf{\% Perdidos}}} & \cellcolor[HTML]{952626}{\textcolor{white}{\textbf{N Categorías}}} & \cellcolor[HTML]{952626}{\textcolor{white}{\textbf{Categoría Más Frecuente}}}\\
\hline
\textbf{Sexo} & 100\% & 0\% & 2 & 1: 1334, 0: 647\\
\hline
\textbf{Nivel socioeconómico} & 100\% & 0\% & 3 & Med: 1466, Alt: 399, Baj: 116\\
\hline
\textbf{Tipo de colegio} & 98.79\% & 1.21\% & 3 & Par: 778, Mun: 683, Par: 496\\
\hline
\end{tabular}}

}

\end{table}%

\subsection{Distinción del tiempo: cohorte y
año}\label{distinciuxf3n-del-tiempo-cohorte-y-auxf1o}

En los datos se encuentran dos variables para medir el tiempo, a saber,
la cohorte y el año, los cuales se distinguen en sus dimensiones
conceptual y analítica, por lo que es fundamental comprender su
diferencia para la interpretación de los análisis.

Las cohortes apuntan al año de ingreso, donde los estudiantes de
cohortes más antiguas tienen mayor información disponible debido a que
han cursado más semestres y, por tanto, más ramos.

El año es una variable que está asociada a cuándo se realizó el curso,
en ese sentido, los cursos cuentan con información que puede someterse a
comparación entregando una interpretación más intuitiva que las
cohortes. Al comparar las notas promedios de los cursos por año, esto
enfoca los cambios que se dan en el tiempo, es decir, entre año y año,
pero tomando en cuenta que los cambios están enfocados en los cursos. En
este caso, las cohortes tienen un rol secundario, ya que personas de
distintas cohortes pueden compartir cursos entre sí, por lo que no
afectarían los resultados.

Los análisis centrados en estudiantes se realizaron utilizando la
variable cohorte para modelar cambios en el tiempo, mientras que los
cambios por año se visualizan con los datos centrados en cursos. Para
descomponer el efecto de la cohorte y el tiempo en las notas, en función
de robustecer el análisis, se generó una sección de
\href{https://tomasurzuam.github.io/estudio-notas/docs/apendices/A-analisis-tiempo.html}{análisis
por año}, la cual se encuentra en anexos.

\subsection{Métodos}\label{muxe9todos}

En una primera etapa se realizan análisis descriptivos de las
calificaciones, considerando medidas de tendencia central y dispersión
(promedio, mediana, desviación estándar y rangos). Estos análisis se
efectúan, en primer lugar, a nivel de Facultad, con el objetivo de
identificar patrones generales. Posteriormente, los datos se desagregan
por carrera y cohorte, permitiendo una comparación más detallada entre
programas académicos y una mejor comprensión de la variabilidad interna
del desempeño estudiantil.

En una segunda etapa se examinan las asociaciones entre las
calificaciones y un conjunto de variables sociodemográficas relevantes.
Las asociaciones se presentan principalmente a partir de gráficos,
siguiendo este orden: nivel facultad (1), nivel carrera (2) y nivel
facultad por cohorte (3). Esta fase busca identificar relaciones que
diluciden la composición de las calificaciones a partir de distintas
variables de interés.

En la tercera etapa se estiman modelos de regresión multivariada con el
fin de evaluar el efecto independiente de cada variable explicativa
sobre las calificaciones, controlando por el resto de los factores
incluidos en el modelo. La especificación de los modelos se ajusta a la
naturaleza continua de la variable dependiente y a los supuestos
estadísticos correspondientes. Se reportan los coeficientes estimados y
niveles de significación, junto con indicadores de ajuste global del
modelo. Como aspecto adicional, se incorporan interacciones entre
variables con un fin exploratorio, con tal de evidenciar si los efectos
de alguna variable por sobre otra son matizados o potenciados.

\section{Sección cursos}\label{secciuxf3n-cursos}

\subsection{Datos}\label{datos-1}

Los datos para la sección de cursos se construyeron a partir de bases de
datos definidas por la cohorte de ingreso de los estudiantes,
realizándose así un proceso de transformación para obtener una base
centrada en cursos. Para esto, se procesaron bases que contienen
información de cursos \textbf{entre el 2018 y el 2025}, tales como el
nombre del curso, su código respectivo y la nota de los y las
estudiantes que estuvieron en tal curso. Al aplicar procesamientos
iniciales, se dio cuenta que la base de 2018 tiene cursos perdidos, es
decir, carece de información en comparación a las bases de los demás
años. Esto junto con información adicional se puede encontrar en los
\href{https://tomasurzuam.github.io/estudio-notas/docs/apendices/C-anexos-cursos.html}{anexos}.

A partir de la información de cursos se logró profundizar en la
operacionalización de la inflación de notas, ya que además de contar con
la calificación promedio de los cursos, se pudieron integrar el
porcentaje de aprobación de los cursos y la dispersión de notas de cada
uno de ellos, logrando la consrtucción de un instrumento de medición
capaz de aportar mayores recursos de información a la comprensión del
concepto y el análisis de los resultados.

Para identificar la carrera de los cursos se realizó un proceso de
clasificación, el cual se basó en \textbf{la identificación del prefijo
del código del curso}. Por ejemplo, si el código del curso empezaba con
``SOC'', se clasificaba como un curso de Sociología. Este proceso se
realizó para cada una de las carreras, utilizando palabras clave
específicas para cada una, lo que nos permitió reconocer la carrera de
origen a la que pertenecen los cursos. En este contexto, el estudio
considera solamente cursos que pertenecen a alguna de las cinco carreras
contempladas (Antropología, Educación Parvularia, Psicología,
Sociología, Trabajo Social), dejando fuera aquellos cursos a nivel
campus tales como artísticos y/o deportivos. Esto marca una distinción
con los cursos disponibles en la sección anterior, donde se consideraron
todos los ramos cursados por estudiante con el objetivo de abarcar la
mayor cantidad de información posible. En cambio, para esta sección, la
decisión se fundamenta en la prioridad que tiene este informe hacia la
Facultad de Ciencias Sociales, enfocándose solamente en las notas que
son asignadas dentro de FACSO.

Para enriquecer el análisis se creó una variable de identificación del
semestre curricular, la cual se calculó restando el año en que los
estudiantes cursaron el ramo con el año de ingreso a la carrera.
Considerando la moda de cada curso, se pudo identificar en qué semestre
curricular se encontraba cada uno de ellos. Con ello, se pudo encontrar
que habían casos disonantes: personas que figuraban con un curso algún
año anterior a su ingreso. Estos casos fueron eliminados debido a su
baja presencia (el año con mayor presencia de casos disonantes fue 2020
con un 3.6\% de los casos totales) y a la falta de información para
hallar su origen.Además, se realizó un proceso de \textbf{clasificación
de los cursos según su carácter}, es decir, una variable que
diferenciaba cursos obligatorios y electivos. Para esto, se utilizó
información de la malla curricular de cada carrera, lo que nos permitió
aplicar una identificación inicial de los cursos obligatorios y
clasificarlos según su área temática. Después de esto, algunos cursos
quedaron sin clasificación, por lo que tuvo que realizarse un proceso
manual para asegurar su correcta categorización.

Terminado la creación de variables, se generó una base madre que
contenía toda la información disponible, y a partir de esta se
desprendieron dos bases agregadas a nivel curso: la primera considerando
las secciones de cada curso, y una segunda que unificaba los cursos que
tuvieran más de una sección.

Finalizado este proceso, se llevó a cabo una depuración final de los
datos, pues habían cursos duplicados y otros que eran inconsistentes.
Los duplicados fueron aquellos que repetían la misma información (nombre
y código del curso, semestre y año), por lo que se optó por mantener
solamente un caso de ellos. Por su parte, los casos inconsistentes eran
cursos que tenían el mismo nombre, año y código pero distinto semestre
curricular. Este problema se generó principalmente en los cursos
electivos debido a que no son cursos con un semestre fijo, pues se
ofrecen a estudiantes entre tercer y quinto año en algunos casos. Para
solucionar esto, se unificaron los cursos en orden de perder la menor
cantidad de información, ponderando sus medidas para minimizar el margen
de error al calcular los valores asociados a las calificaciones.

Las últimas decisiones de análisis se cristalizan en dos principalmente:
el uso de todos los cursos disponibles con sus respectivas secciones en
el análisis (1) y la consideración de cursos con, al menos, 10
estudiantes (2). La primera decisión se justifica en que cada sección de
un mismo curso se puede comprender como un curso en sí mismo, lo que
constituye un aporte valioso para la variabilidad de información, así
como para observar de manera más clara el panorama. La segunda decisión
se fundamenta en la alta influencia que tenían los cursos con pocos
estudiantes en los análisis de resultados. En este marco, la mayoría de
los cursos de tamaño reducido tenían solamente un estudiante, los cuales
correspondían a cursos de seminario y/o práctica. Al estar registrados
como cursos, la evaluación a una persona tenía el mismo peso que una
evaluación a 100 estudiantes, lo que no es acertado en términos
estadísticos ni sustantivos. Para justificar esto empíricamente, el
promedio general de la facultad con todos los datos es de un 6.4, y al
dejar fuera los cursos con menos de 10 personas disminuye a un 6.0 (ver
\href{https://tomasurzuam.github.io/estudio-notas/docs/apendices/C-anexos-cursos.html}{anexos}),
demostrando la distorsión de los resultados sin el filtrado. Por ello,
se tomó la decisión de conservar para el análisis aquellos cursos con un
mínimo de 10 personas.

A partir de todas las decisiones de procesamiento tomadas, la muestra
analítica se constituyó en un total de 2291 cursos.

\subsection{Variables}\label{variables-1}

En la Tabla~\ref{tbl-codebook-cursos} se pueden visualizar todas las
variables disponibles para el análisis de cursos junto con su naturaleza
y una breve descripción de cada una. Posteriormente, se presentan los
descriptivos de las variables, distinguiendo por variables categóricas
(Tabla~\ref{tbl-categoricas-cursos}) y variables numéricas
(Tabla~\ref{tbl-numericas-cursos}).

\begin{table}

\caption{\label{tbl-codebook-cursos}Libro de códigos}

\centering{

\begingroup\fontsize{9}{11}\selectfont

\resizebox{\ifdim\width>\linewidth\linewidth\else\width\fi}{!}{
\begin{tabu} to \linewidth {>{\raggedright}X>{\raggedright}X>{\centering}X>{\raggedright}X}
\hline
\cellcolor[HTML]{952626}{\textcolor{white}{\textbf{Grupo}}} & \cellcolor[HTML]{952626}{\textcolor{white}{\textbf{Variable}}} & \cellcolor[HTML]{952626}{\textcolor{white}{\textbf{Tipo}}} & \cellcolor[HTML]{952626}{\textcolor{white}{\textbf{Descripción}}}\\
\hline
 & \ttfamily{\textbf{codigo}} & Categórica & Código único del curso (ej. ANT00001). Primeros 3 caracteres = sigla carrera; 5 siguientes = correlativo.\\
\cline{2-4}
 & \ttfamily{\textbf{curso}} & Categórica & Nombre oficial del curso según catálogo académico.\\
\cline{2-4}
 & \ttfamily{\textbf{carrera}} & Categórica & Nombre del programa académico al que pertenece el curso.\\
\cline{2-4}
\multirow[t]{-4}{2.2cm}{\raggedright\arraybackslash \textbf{Identificación}} & \ttfamily{\textbf{caracter}} & Categórica & Carácter del curso en el plan de estudios: 'Obligatorio' o 'Electivo'.\\
\cline{1-4}
 & \ttfamily{\textbf{ano}} & Categórica & Año académico en que se dictó la sección. Convertir con as.integer(ano) para cálculos.\\
\cline{2-4}
 & \ttfamily{\textbf{semestre\_moda}} & Numérica & Semestre modal del curso según la malla curricular. Rango: 1–10.\\
\cline{2-4}
 & \ttfamily{\textbf{semestre\_cv}} & factor ordenado & Semestre del plan de estudios como etiqueta textual. 10 niveles: Primer semestre … Décimo semestre.\\
\cline{2-4}
\multirow[t]{-4}{2.2cm}{\raggedright\arraybackslash \textbf{Temporal}} & \ttfamily{\textbf{seccion}} & integer & Número de sección del curso en el año. Un mismo curso puede tener varias secciones por año.\\
\cline{1-4}
 & \ttfamily{\textbf{n\_estudiantes}} & integer & Total de estudiantes matriculados al cierre del período (excluye bajas oficiales).\\
\cline{2-4}
 & \ttfamily{\textbf{n\_aprobados}} & integer & Estudiantes que aprobaron el curso (nota ≥ 4.0). Verificar: n\_aprobados + n\_reprobados = n\_estudiantes.\\
\cline{2-4}
 & \ttfamily{\textbf{n\_reprobados}} & integer & Estudiantes que reprobaron el curso (nota < 4.0).\\
\cline{2-4}
 & \ttfamily{\textbf{n\_hombres}} & integer & Estudiantes hombres matriculados según registro oficial. Verificar: n\_hombres + n\_mujeres = n\_estudiantes.\\
\cline{2-4}
\multirow[t]{-5}{2.2cm}{\raggedright\arraybackslash \textbf{Matrícula}} & \ttfamily{\textbf{n\_mujeres}} & integer & Estudiantes mujeres matriculadas según registro oficial.\\
\cline{1-4}
 & \ttfamily{\textbf{nota\_promedio}} & Numérica & Nota promedio de los estudiantes de la sección. Escala 1.0–7.0; aprobación mínima: 4.0.\\
\cline{2-4}
 & \ttfamily{\textbf{desv\_std}} & Numérica & Desviación estándar de las notas de la sección. Indica dispersión del rendimiento.\\
\cline{2-4}
 & \ttfamily{\textbf{nota\_min}} & Numérica & Nota mínima obtenida en la sección. Escala 1.0–7.0.\\
\cline{2-4}
 & \ttfamily{\textbf{nota\_max}} & Numérica & Nota máxima obtenida en la sección. Escala 1.0–7.0.\\
\cline{2-4}
\multirow[t]{-5}{2.2cm}{\raggedright\arraybackslash \textbf{Rendimiento}} & \ttfamily{\textbf{porcentaje\_aprobacion}} & Numérica & Porcentaje de aprobación: (n\_aprobados / n\_estudiantes) × 100. Rango: 0–100.\\
\hline
\end{tabu}}
\endgroup{}

}

\end{table}%

\begin{table}

\caption{\label{tbl-categoricas-cursos}Estadísticas descriptivas de
variables categóricas}

\centering{

\centering\begingroup\fontsize{8}{10}\selectfont

\begin{tabular}{>{\raggedright\arraybackslash}p{2.2cm}|>{\centering\arraybackslash}p{1.8cm}|>{\centering\arraybackslash}p{1.7cm}|>{\centering\arraybackslash}p{1.7cm}|>{\raggedright\arraybackslash}p{5cm}}
\hline
\multicolumn{1}{c|}{\cellcolor[HTML]{952626}{\textcolor{white}{\textbf{ }}}} & \multicolumn{2}{c|}{\cellcolor[HTML]{952626}{\textcolor{white}{\textbf{Muestra}}}} & \multicolumn{2}{c}{\cellcolor[HTML]{952626}{\textcolor{white}{\textbf{Categorías}}}} \\
\cline{2-3} \cline{4-5}
\cellcolor[HTML]{952626}{\textcolor{white}{\textbf{Variable}}} & \cellcolor[HTML]{952626}{\textcolor{white}{\textbf{\% Casos válidos}}} & \cellcolor[HTML]{952626}{\textcolor{white}{\textbf{\% Perdidos}}} & \cellcolor[HTML]{952626}{\textcolor{white}{\textbf{N Categorías}}} & \cellcolor[HTML]{952626}{\textcolor{white}{\textbf{Distribución}}}\\
\hline
\textbf{Semestre curricular} & 99.91\% & 0.09\% & 10 & Sexto semestre (n = 425); Séptimo semestre (n = 306); Quinto semestre (n = 292); Cuarto semestre (n = 271); Segundo semestre (n = 254); Tercer semestre (n = 247); Primer semestre (n = 243); Octavo semestre (n = 225); Noveno semestre (n = 16); Décimo semestre (n = 10)\\
\hline
\textbf{Tipo de curso} & 98.6\% & 1.4\% & 2 & Obligatorio (n = 1616); Electivo (n = 643)\\
\hline
\textbf{Carrera} & 100\% & 0\% & 5 & Antropología (n = 610); Psicología (n = 602); Sociología (n = 510); Educación Parvularia (n = 289); Trabajo Social (n = 280)\\
\hline
\end{tabular}
\endgroup{}

}

\end{table}%

\begin{table}

\caption{\label{tbl-numericas-cursos}Estadísticas descriptivas de
variables numéricas}

\centering{

\centering\begingroup\fontsize{13}{15}\selectfont

\resizebox{\ifdim\width>\linewidth\linewidth\else\width\fi}{!}{
\begin{tabular}{>{\raggedright\arraybackslash}p{12em}|c|c|c|c|c|c|c|c|c}
\hline
\multicolumn{1}{c|}{\cellcolor[HTML]{952626}{\textcolor{white}{\textbf{ }}}} & \multicolumn{2}{c|}{\cellcolor[HTML]{952626}{\textcolor{white}{\textbf{Muestra}}}} & \multicolumn{7}{c}{\cellcolor[HTML]{952626}{\textcolor{white}{\textbf{Estadísticos descriptivos}}}} \\
\cline{2-3} \cline{4-10}
\cellcolor[HTML]{952626}{\textcolor{white}{\textbf{Variable}}} & \cellcolor[HTML]{952626}{\textcolor{white}{\textbf{\% Casos válidos}}} & \cellcolor[HTML]{952626}{\textcolor{white}{\textbf{\% Perdidos}}} & \cellcolor[HTML]{952626}{\textcolor{white}{\textbf{Media}}} & \cellcolor[HTML]{952626}{\textcolor{white}{\textbf{DE}}} & \cellcolor[HTML]{952626}{\textcolor{white}{\textbf{Mín}}} & \cellcolor[HTML]{952626}{\textcolor{white}{\textbf{Q1}}} & \cellcolor[HTML]{952626}{\textcolor{white}{\textbf{Mediana}}} & \cellcolor[HTML]{952626}{\textcolor{white}{\textbf{Q3}}} & \cellcolor[HTML]{952626}{\textcolor{white}{\textbf{Máx}}}\\
\hline
\textbf{Nota promedio} & 100\% & 0\% & 6.01 & 0.52 & 3.76 & 5.70 & 6.07 & 6.39 & 7.00\\
\hline
\textbf{Desviación estándar} & 100\% & 0\% & 0.67 & 0.38 & 0.00 & 0.39 & 0.59 & 0.88 & 2.67\\
\hline
\textbf{Cantidad de estudiantes} & 100\% & 0\% & 42.42 & 25.87 & 10.00 & 24.00 & 38.00 & 54.00 & 197.00\\
\hline
\textbf{Tasa de aprobación} & 100\% & 0\% & 97.28 & 4.69 & 55.81 & 96.15 & 100.00 & 100.00 & 100.00\\
\hline
\textbf{Cantidad de hombres} & 100\% & 0\% & 12.98 & 10.31 & 0.00 & 5.00 & 12.00 & 18.00 & 71.00\\
\hline
\textbf{Cantidad de mujeres} & 100\% & 0\% & 29.44 & 18.43 & 1.00 & 15.00 & 26.00 & 39.00 & 126.00\\
\hline
\end{tabular}}
\endgroup{}

}

\end{table}%

\subsection{Métodos}\label{muxe9todos-1}

La primera parte del análisis centrado en cursos considera la
presentación general de estadísticos descriptivos con el objetivo de
entregar una comprensión profunda y clara de la composición de los
datos. Se comienza con descriptivos relacionados a la cantidad de cursos
existentes en FACSO, para luego dar paso a los descriptivos que dan
cuenta de distintas maneras la composición de las calificaciones en la
facultad: la nota, el porcentaje de aprobación y la dispersión de las
notas. En esta sección se presentan gráficos a nivel facultad y
desagregados por carrera, primero desde una perspectiva general y
posteriormente presentando las trayectorias en el tiempo de las
variables.

En la segunda parte se presentan los cruces entre las notas de los
cursos y las variables de caracterización: tipo de curso, género y
tamaño del curso. Las asociaciones se ordenan a partir de la siguiente
estructura de gráficos: gráfico general de facultad (1), trayectoria
temporal de la facultad (2), gráfico desagregado por carreras (3) y
trayectoria temporal desagregada por carreras (4). En este punto se
pretende identificar en qué medida las variables de caracterización
determinan las calificaciones de los cursos, además de buscar distinguir
patrones generalizables a nivel facultad tanto como hallazgos
particulares de carreras.

\bookmarksetup{startatroot}

\chapter{Resultados: Estudiantes}\label{resultados-estudiantes}

Este apartado tiene como objetivo presentar los principales resultados
de la investigación con los estudiantes como unidad de análisis. En
primer lugar, se exponen resultados descriptivos a nivel de Facultad y,
posteriormente, se detallan a nivel de carrera. En segundo lugar, se
analizan las asociaciones entre las notas y las variables de
caracterización más relevantes (sexo, colegio de procedencia, nivel
socioeconómico y prueba de admisión universitaria). Por último, se
estima un modelo de regresión con el objetivo de identificar qué
variables resultan más determinantes en las calificaciones de las y los
estudiantes.

La presentación de los resultados descriptivos y de asociación se
organiza de la siguiente manera: (1) tabla de promedios de notas por
carrera y categoría; (2) gráfico de distribución de las notas a nivel de
Facultad y (3) gráficos de distribución de las notas según carrera.

\section{Descriptivos generales}\label{descriptivos-generales}

Esta sección tiene como objetivo determinar el promedio general de los
estudiantes de FACSO. En primer lugar, se va a estimar la distribución
de las notas a partir de la base de promedios generales UCampus, la cual
tiene un N total de 2578.

\begin{figure}

\centering{

\pandocbounded{\includegraphics[keepaspectratio]{03-resultados-est_files/figure-pdf/fig-histo-og-1.pdf}}

}

\caption{\label{fig-histo-og}Elaboración propia - datos de promedios
generales UCampus (2026)}

\end{figure}%

Considerando la Figura~\ref{fig-histo-og}, se confirma que la estimación
del promedio general realizada en el análisis previo por la Dirección de
Pregrado de la facultad es correcta. El gráfico presenta una
distribución asimétrica, donde la mayoría de los casos se concentran en
notas entre 6.0 y 7.0, mientras que las notas menores a 5.0 son
inperceptibles, destacándose un promedio a nivel facultad de 6.4.
Además, el peak de concentración de notas está por sobre el promedio, es
decir, cercano al 6.6-6.7. Respecto a la dispersión de las notas, el
promedio de la base original general cuenta con una desviación estándar
de 0.33, lo que significa que aproximadamente dos tercios de los
estudiantes tienen un promedio final de notas entre de 6.1 y 6.7
(asumiendo distribución normal). En suma, la información entregada por
la figura refleja una concentración de las notas de los estudiantes en
la parte más alta de la escala de notas, acompañado de una baja
dispersión.

Ahora bien, con la base general de promedios de UCampus no teníamos
información clara respecto al proceso de construcción del promedio
general de cada estudiante. Debido a ello, se solicitó una base con las
notas desagregadas por curso para así calcular el promedio de
estudiantes a partir de sus notas en cada curso. Al realizar este
cálculo, el resultado es el siguiente:

\begin{figure}

\centering{

\pandocbounded{\includegraphics[keepaspectratio]{03-resultados-est_files/figure-pdf/fig-histo-c-1.pdf}}

}

\caption{\label{fig-histo-c}Elaboración propia - datos de notas por
cursos UCampus (2026)}

\end{figure}%

Se observa que la Figura~\ref{fig-histo-c} tiene una distribución de las
calificaciones menos concentrada en la parte alta de la escala de notas,
lo cual se condice con una disminución del promedio a nivel facultad a
un 6.1. En relación con esto, la desviación estándar corresponde a un
0.62, lo que significa que el 68\% de los casos (asumiendo distribución
normal) se encuentra entre un 5.5 y un 6.7, afirmando que este cálculo
del promedio presenta una mayor dispersión en las notas. Aun así, el
promedio general sigue siendo alto e, incluso, el 50\% de los casos se
concentra por encima del 6.2.

La disminución del promedio respecto a la Figura~\ref{fig-histo-og} se
debe a las decisiones tomadas para el cálculo de los promedios. En este
caso, se decidió considerar las notas de los cursos aprobados y
reprobados por cada estudiante para estimar su promedio final. Esto se
debe al enfoque del estudio, puesto que nos interesa reflejar de mejor
manera el proceso de calificación en la facultad más que solamente la
nota final. Otra razón para considerar este promedio de notas es que
este cálculo puede reflejar mayor varianza y, por tanto, entregar más
información respecto a la asociación con variables de caracterización.
Sin embargo, las variables de caracterización se encontraban solo en la
base inicial de promedios generales, por lo que se realizó una fusión de
la base original general (que cuenta con variables de caracterización) y
la base de notas por curso. Esto trajo el siguiente resultado:

\begin{figure}

\centering{

\pandocbounded{\includegraphics[keepaspectratio]{03-resultados-est_files/figure-pdf/fig-histo-facso-1.pdf}}

}

\caption{\label{fig-histo-facso}Elaboración propia en base a datos de
UCampus (2026)}

\end{figure}%

Al unir las dos bases se reducen los casos a 1981 estudiantes. La
primera diferencia que se observa en la Figura~\ref{fig-histo-facso} es
que el promedio disminuye a un 5.9. Esto se complementa con la
distribución de las notas, donde la columna que representa las
calificaciones más altas ya no es la que tiene la mayor concentración.
La desviación estándar corresponde a 0.6, lo que significa que el 68\%
de los casos se mantiene entre el 5.3 y el 6.5, presentando una mayor
variabilidad que la Figura~\ref{fig-histo-og}. Cabe considerar que este
es el escenario con el resultado más bajo que se puede encontrar. Aun
así, la mitad de los estudiantes obtienen notas sobre 6.0, incluyendo
los cursos reprobados en sus trayectorias. En síntesis, si bien este
cálculo del promedio de los estudiantes permite observar una mayor
dispersión respecto al promedio que contenía la base original general,
la concentración de las calificaciones sigue ubicándose en la parte
superior de la escala de notas.

\begin{figure}

\centering{

\pandocbounded{\includegraphics[keepaspectratio]{03-resultados-est_files/figure-pdf/fig-box-facso-1.pdf}}

}

\caption{\label{fig-box-facso}Elaboración propia en base a datos de
UCampus (2026)}

\end{figure}%

Para profundizar en la comprensión del promedio general de la facultad,
en la Figura~\ref{fig-box-facso} se presentan las diferencias en los
promedios de notas según carrera. Se observa que Educación Parvularia
presenta los valores centrales más, seguido de Psicolgía. Antropología,
Sociología y Trabajo Social exhiben valores muy similares. Cabe resaltar
que en todas las carreras se observan algunos valores atípicos bajos, lo
que indica la existencia de muy pocos estudiantes con promedios
considerablemente inferiores al promedio general.

\section{Asociación entre nota y variables de
caracterización}\label{asociaciuxf3n-entre-nota-y-variables-de-caracterizaciuxf3n}

\subsection{Género}\label{guxe9nero}

\begin{table}

\centering{

\centering\begingroup\fontsize{14}{16}\selectfont

\resizebox{\ifdim\width>\linewidth\linewidth\else\width\fi}{!}{
\begin{tabular}{l|c|c|c}
\hline
Carrera & Hombre & Mujer & Diferencia (Mujer - Hombre)\\
\hline
Antropología & 5.79 & 5.93 & 0.14\\
\hline
Educación Parvularia & 5.90 & 6.20 & 0.30\\
\hline
Psicología & 5.86 & 6.07 & 0.20\\
\hline
Sociología & 5.80 & 5.92 & 0.12\\
\hline
Trabajo Social & 5.41 & 5.89 & 0.48\\
\hline
\end{tabular}}
\endgroup{}

}

\caption{\label{tbl-gen}Promedios por carrera según género 2020-2024}

\end{table}%

La Tabla~\ref{tbl-gen} exhibe los promedios de notas por carrera
distinguiendo entre hombres y mujeres. Se observa una tendencia
consistente en las cinco carreras del estudio, donde las mujeres
sistemáticamente obtienen mejores calificaciones. En este contexto,
Trabajo Social es la carrera que presenta la mayor brecha en rendimiento
académico por género, con 0.48 décimas a favor de las estudiantes.

\begin{figure}

\centering{

\pandocbounded{\includegraphics[keepaspectratio]{03-resultados-est_files/figure-pdf/fig-hist-gen-1.pdf}}

}

\caption{\label{fig-hist-gen}Elaboración propia en base a datos de
UCampus (2026)}

\end{figure}%

En la Figura~\ref{fig-hist-gen} se encuentran histogramas que presentan
la distribución de las notas por género, primero a nivel facultad y
luego por carrera. Un primer hallazgo a reportar es la forma de la
distribución entre hombres y mujeres, las cuales son claramente
distintas. Aquello se evidencia en que los hombres presentan más casos
en la cola izquierda de la distribución, aludiendo a notas más bajas,
mientras que las mujeres tienen una concentración mayor desde el 6 en
adelante. Trabajo Social demuestra claramente este resultado, donde
destacan varias barras pequeñas azules entre el 1 y el 5, contrario a la
menor dispersión de las barras naranjas, acumuladas entre el 6 y 6.3.
Por último, en Educación Parvularia no se logran divisar barras de
hombres, lo que se debe a la mínima cantidad de estudiantes varones que
cursan esa carrera.

\subsection{Tipo de colegio}\label{tipo-de-colegio}

\begin{table}

\centering{

\centering\begingroup\fontsize{14}{16}\selectfont

\resizebox{\ifdim\width>\linewidth\linewidth\else\width\fi}{!}{
\begin{tabular}{l|c|c|c}
\hline
carrera & Municipal & Particular Subvencionado & Particular Pagado\\
\hline
Antropología & 5.79 & 5.89 & 5.93\\
\hline
Educación Parvularia & 6.12 & 6.26 & 6.27\\
\hline
Psicología & 5.79 & 6.16 & 6.14\\
\hline
Sociología & 5.73 & 5.87 & 5.96\\
\hline
Trabajo Social & 5.72 & 5.84 & 5.60\\
\hline
\end{tabular}}
\endgroup{}

}

\caption{\label{tbl-notas-col}Promedios por carrera según colegio
2020-2024}

\end{table}%

La Tabla~\ref{tbl-notas-col} muestra los promedios por carrera,
distinguiendo el tipo de colegio del que provienen los estudiantes. En
antropología, educación parvularia y sociología se aprecia un patrón
escalonado ascendente según el financiamiento de las escuelas, donde las
notas más bajas son de aquellos que egresaron de un colegio municipal,
siendo superados por los particulares subvencionados, y los que
estudiaron en particulares pagados son quienes poseen mejores
calificaciones. Sin embargo, en Psicología y Trabajo Social se presenta
un patrón distinto al ya visto, pues el grupo de particulares
subvencionados es el que obtiene las calificaciones más altas, mientras
que los particulares pagados son relegados al segundo lugar.

\begin{figure}

\centering{

\pandocbounded{\includegraphics[keepaspectratio]{03-resultados-est_files/figure-pdf/fig-dens-carreras-1.pdf}}

}

\caption{\label{fig-dens-carreras}Elaboración propia en base a datos de
UCampus (2026)}

\end{figure}%

La Figura~\ref{fig-dens-carreras} presenta las distribuciones de
densidad, primero en un nivel de facultad, y posteriormente por cada
carrera. A nivel facultad se observa claramente que los colegios
municipales poseen una concentración menor de notas desde el 6 en
adelante, a comparación con los particulares subvencionados y pagados.
Las curvas verdes y azules mantienen una forma bastante similar, lo que
da a entender que, en promedio en FACSO, quienes provienen de algún tipo
de colegio particular, tienen un rendimiento académico similar. En
sociología y psicología se mantiene este patrón, no así en antropología
y educación parvularia, donde los particulares subvencionados tienen una
concentración mayor entre el 6 y 6.5 por sobre los particulares pagados,
mientras los municipales se mantienen en último lugar. Sin embargo, en
trabajo social se da algo distinto a las demás carreras, y es que las
notas de los provenientes de particulares pagados son las más bajas, lo
que queda demostrado al ser la curva más baja de las tres, seguido de
los municipales y los subvencionados.

\subsection{Nivel socioeconómico}\label{nivel-socioeconuxf3mico}

La variable nivel socioeconómico fue obtenida mediante una encuesta
institucional de caracterización aplicada de manera externa a la
plataforma UCampus, siendo posteriormente integrada a sus bases de
datos. Esta variable refleja el nivel de vulnerabilidad de la comuna de
origen del estudiante, con valores que oscilan entre 0 y 8, donde los
valores más bajos indican mayor vulnerabilidad socioeconómica. Con el
propósito de facilitar su interpretación y análisis, la variable fue
recategorizada en tres niveles: bajo (0, 1 y 2), medio (3, 4 y 5) y alto
(6, 7 y 8).

\begin{table}

\centering{

\centering\begingroup\fontsize{14}{16}\selectfont

\resizebox{\ifdim\width>\linewidth\linewidth\else\width\fi}{!}{
\begin{tabular}{l|c|c|c}
\hline
carrera & Bajo & Medio & Alto\\
\hline
Antropología & 5.76 & 5.85 & 6.01\\
\hline
Educación Parvularia & 5.91 & 6.25 & 6.17\\
\hline
Psicología & 5.94 & 5.99 & 6.07\\
\hline
Sociología & 5.65 & 5.85 & 5.93\\
\hline
Trabajo Social & 5.79 & 5.77 & 5.65\\
\hline
\end{tabular}}
\endgroup{}

}

\caption{\label{tbl-notas-nse}Promedios por carrera según nivel
socioeconómico 2020-2024}

\end{table}%

La Tabla~\ref{tbl-notas-nse} presenta los promedios de nota de las
carreras, distinguiendo por nivel socioeconómico. En los casos de
antropología, psicología y sociología se exhibe un patrón ascendente de
notas desde el menor al mayor estatus socioeconómico, reflejando que,
aquellos con mejores condiciones económicas de base obtienen un logro
académico superior. No obstante, es interesante observar que esta
tendencia no se cumple en educación parvularia ni trabajo social.
Respecto a educación parvularia, son los estudiantes de nivel
socioeconómico medio quienes tienen calificaciones más altas que
aquellos con estatus más alto, mientras que los de nivel bajo siguen
siendo los de menor rendimiento. Por su parte, en trabajo social, los
estudiantes de nivel socioeconómico más alto son los que tienen las
notas más bajas, al contrario de los niveles bajos y medios que,
coincidentemente, poseen la misma nota promediada.

\begin{figure}

\centering{

\pandocbounded{\includegraphics[keepaspectratio]{03-resultados-est_files/figure-pdf/fig-dens-carreranse-1.pdf}}

}

\caption{\label{fig-dens-carreranse}Elaboración propia en base a datos
de UCampus (2026)}

\end{figure}%

La Figura~\ref{fig-dens-carreranse} presenta gráficos de densidad del
promedio de notas según nivel socioeconómico (NSE), primero a nivel de
facultad y luego desagregados por carrera. A nivel agregado se observa
un patrón estratificado claro: el NSE bajo presenta una menor
concentración de densidad en notas superiores a 6, seguido del NSE
medio, mientras que el NSE alto concentra su distribución en los valores
más altos del rendimiento académico.

En Sociología se observa un patrón distinto, ya que el NSE bajo exhibe
una distribución bimodal, con una mayor frecuencia cercana al 6
---aunque levemente desplazado hacia la izquierda--- y otra en torno a
la nota 5. El NSE alto, en cambio, concentra sus valores hacia el
extremo derecho de la distribución, pese a presentar menor densidad. Por
su parte, psicología reproduce un patrón estratificado similar al de la
facultad.

En Antropología y Educación Parvularia destaca que el NSE medio presenta
una mayor concentración en torno a la nota 6 en comparación con los
otros grupos; en Antropología, además, la distribución del NSE bajo
sugiere una estructura multimodal, lo que indica mayor heterogeneidad en
el desempeño. Finalmente, en Trabajo Social el NSE medio concentra la
mayor densidad cerca del 6, mientras que el NSE bajo presenta
adicionalmente una masa relevante en torno a notas cercanas a 4, y el
NSE alto muestra una distribución más extendida entre aproximadamente
4.5 y 6.

\subsection{Asociación de nota y prueba de admisión
(PSU/PAES)}\label{asociaciuxf3n-de-nota-y-prueba-de-admisiuxf3n-psupaes}

\begin{figure}

\centering{

\pandocbounded{\includegraphics[keepaspectratio]{03-resultados-est_files/figure-pdf/fig-scat-facso-1.pdf}}

}

\caption{\label{fig-scat-facso}Elaboración propia en base a datos de
UCampus (2026)}

\end{figure}%

En la Figura~\ref{fig-scat-facso} se presenta la asociación entre el
promedio de notas de los estudiantes y su percentil en la prueba de
admisión universitaria (PSU o PAES). A nivel facultad, se encontró una
relación positiva pero débil (\emph{r} = .22), sin haber un patrón claro
en cómo se distribuyen los puntos en el plano. Además se observa que la
dispersión es levemente más pronunciada en los percentiles bajos y
medios, donde coexisten estudiantes con notas muy disímiles, mientras
que en los percentiles altos la distribución tiende a concentrarse,
aunque las diferencias son poco notorias. Por lo tanto, el desempeño
académico en la universidad no parece estar relacionado con el desempeño
previo.

\section{Modelamiento}\label{modelamiento}

\begin{figure}

\centering{

\pandocbounded{\includegraphics[keepaspectratio]{03-resultados-est_files/figure-pdf/fig-coefplot-1.pdf}}

}

\caption{\label{fig-coefplot}Elaboración propia en base a datos de
UCampus (2026)}

\end{figure}%

En la Figura~\ref{fig-coefplot} se presentan los coeficientes de
regresión de un modelo multivariado con los datos a nivel de facultad,
donde el promedio de notas es la variable predicha. En primer lugar, las
mujeres obtienen, en promedio, 0.208 puntos más que los hombres (β =
0.208***), diferencia estadísticamente significativa y consistente con
los resultados bivariados. Este hallazgo confirma una ventaja
sistemática femenina en el desempeño académico dentro de la facultad,
independientemente de las demás variables del modelo.

En segundo lugar, los resultados muestran que tanto los estudiantes de
nivel socioeconómico medio (β = 0.145) como alto (β = 0.176) obtienen
promedios significativamente superiores a aquellos estudiante de estrato
bajo. La magnitud de ambos coeficientes es similar, aunque el nivel alto
presenta una ventaja leve.

En tercer lugar, tomando como referencia el colegio municipal, tanto los
estudiantes provenientes de establecimientos subvencionados (β = 0.126)
como particulares pagados (β = 0.092) presentan un rendimiento
significativamente superior. Resulta llamativo que la ventaja de los
subvencionados sobre el municipal sea mayor que la de los particulares
pagados, lo que podría sugerir que, una vez controladas otras variables
como el nivel socioeconómico, el origen en un establecimiento particular
pagado no representa una ventaja adicional tan marcada como podría
esperarse.

Por último, al considerar la carrera, Educación Parvularia destaca con
el coeficiente más alto de todo el modelo (β = 0.416), lo que indica que
sus estudiantes obtienen promedios notablemente superiores a los de
Psicología, confirmando que esta carrera concentra el rendimiento
académico más elevado de la facultad. Por su parte, Trabajo Social (β =
−0.006) y Sociología (β = 0.005) muestran coeficientes prácticamente
nulos y no significativos, lo que indica que su rendimiento es
estadísticamente indistinguible del de Psicología. Antropología presenta
el único coeficiente negativo y marginalmente significativo (β =
−0.067), sugiriendo un rendimiento levemente inferior al de Psicología.

\bookmarksetup{startatroot}

\chapter{Resultados: Cursos}\label{resultados-cursos}

Este apartado tiene como objetivo presentar los principales resultados
de la investigación con los cursos como unidad de análisis. En primer
lugar, se exponen resultados descriptivos de la cantidad de cursos
disponibles para introducir un contexto general, y luego se presentan
los descriptivos de tres variables que son una expresión relevante de
las calificaciones: las notas, el porcentaje de aprobación y la
dispersión de las notas.

Posterior a los descriptivos, se presentan análisis de cruces entre las
notas y tres variables de caracterización de los cursos: el tipo de
curso, la composición de género y el tamaño de los cursos.

La presentación de los resultados se estructura de la siguiente manera:
gráficos descriptivos de facultad (1), trayectoria temporal de la
facultad (2), análisis desagregado por carreras (3) y trayectoria
temporal desagregada por carreras (4). Esta estructura presenta una
excepción en el punto \textbf{5.3 Tasa de aprobación}, donde se muestran
tablas en vez de gráficos para facilitar la interpretación.

\section{Descriptivos generales}\label{descriptivos-generales-1}

Antes de entrar de lleno al análisis de las notas y sus aristas
asociadas, es importante comprender la estructura de cursos de la
Facultad de Ciencias Sociales. A modo de introducción, se presentan
algunos descriptivos clave para entender el el desarrollo del informe,
así como las decisiones más relevantes a hora de qué cursos considerar
en el análisis.

\subsection{Contexto de la conformación de
cursos}\label{contexto-de-la-conformaciuxf3n-de-cursos}

La cantidad de cursos que se imparten en la Facultad de Ciencias
Sociales no es un dato que tenga un registro claro, y sin embargo es una
información valiosa a la hora de estudiar cómo se distribuyen las
calificaciones.

\subsubsection{Total de cursos}

\begin{figure}

\centering{

\pandocbounded{\includegraphics[keepaspectratio]{04-resultados-curso_files/figure-pdf/fig-cursos-facso-1.pdf}}

}

\caption{\label{fig-cursos-facso}Elaboración propia (2026)}

\end{figure}%

\subsubsection{Distinción por curso}

\begin{figure}

\centering{

\pandocbounded{\includegraphics[keepaspectratio]{04-resultados-curso_files/figure-pdf/fig-cursos-facso-stack-1.pdf}}

}

\caption{\label{fig-cursos-facso-stack}Elaboración propia (2026)}

\end{figure}%

En la Figura~\ref{fig-cursos-facso} se observa la cantidad de cursos por
año que ofrece la facultad en Sociología, Psicología, Antropología,
Trabajo Social y Educación Parvularia. El resultado más llamativo y
claro es el patrón ascendente en la cantidad de cursos conforme el paso
del tiempo. Si bien el crecimiento de los cursos parece paulatino por la
forma de las barras, resulta interesante que los cursos en el año 2024
crecen de manera exponencial, marcando una diferencia de más de 300
cursos con el año anterior. En la Figura~\ref{fig-cursos-facso-stack} se
muestra el mismo análisis pero desagregado por los cursos con sección
única y aquellos con más de una sección. Aquí hay una pista clara: el
crecimiento de los cursos se debe principalmente al incremento de cursos
multisección. Entre 2021 y 2025 los cursos de sección única se mantienen
estables, pero son aquellos con más de una sección los que terminan por
marcar una pauta clara en la oferta programática. Esto puede también
deberse a la formalización de cursos en las plataformas de la
universidad que, en años anteriores, no contaban con un registro
oficial.

\begin{figure}

\centering{

\pandocbounded{\includegraphics[keepaspectratio]{04-resultados-curso_files/figure-pdf/fig-cursos-carrera-stack-1.pdf}}

}

\caption{\label{fig-cursos-carrera-stack}Elaboración propia (2026)}

\end{figure}%

En la Figura~\ref{fig-cursos-carrera-stack} se presentan los cursos que
existen por carrera entre el 2018 y el 2025. Se pueden observar
diferencias claras entre carreras: Antropología y Psicología son las
carreras que ofrecen más cursos en la facultad. Por su parte, Sociología
ofrece menos cursos que las dos carreras anteriores, pero se caracteriza
por tener la mayor cantidad de cursos con más de una sección en términos
proporcionales como absolutos. Educación Parvularia y Trabajo Social son
las carreras que tienen la oferta programática más baja, ofreciendo
menos de la mitad de cursos que las tres carreras ya mencionadas.

\begin{figure}

\centering{

\pandocbounded{\includegraphics[keepaspectratio]{04-resultados-curso_files/figure-pdf/fig-facso-notasnof-1.pdf}}

}

\caption{\label{fig-facso-notasnof}Elaboración propia (2026)}

\end{figure}%

La Figura~\ref{fig-facso-notasnof} muestra la distribución de notas en
FACSO entre 2018 y 2025, la cual es marcadamente asimétrica, con una
fuerte concentración de notas entre 6 y 7 y donde el valor más frecuente
es precisamente 7 (cerca de 500 casos). La línea punteada vertical
representa el promedio, ubicándose aproximadamente en 6.3 al considerar
el curso como unidad de análisis. Las notas bajo 5 son muy poco
frecuentes, y prácticamente no existen registros entre 2 y 3.

\begin{figure}

\centering{

\pandocbounded{\includegraphics[keepaspectratio]{04-resultados-curso_files/figure-pdf/fig-density-est-1.pdf}}

}

\caption{\label{fig-density-est}Elaboración propia (2026)}

\end{figure}%

La Figura~\ref{fig-density-est} muestra la concentración de estudiantes
por cursos a nivel facultad. El gráfico entrega un resultado importante:
el volumen mayoritario de cursos tiene aproximadamente entre 1 a 10
estudiantes. Si bien hay una parte que cubre hasta los 60 estudiantes
aproximadamente, la curva de la parte izquierda de la figura es
sumamente pronunciada, resaltando con creces en comparación a las otras
áreas. Este resultado es llamativo, pero es necesario profundizar un
poco más en el análisis para encontrar alguna razón a ello.

\begin{figure}

\centering{

\pandocbounded{\includegraphics[keepaspectratio]{04-resultados-curso_files/figure-pdf/fig-semestres-1.pdf}}

}

\caption{\label{fig-semestres}Elaboración propia (2026)}

\end{figure}%

Ante este panorama, es necesario ahondar en la composición de los cursos
para entender en qué medida las calificaciones de distintos cursos son
comparables. Al respecto, un elemento fundamental es el número de
estudiantes por curso, ya que carece de sentido comparar cursos muy
distintos en esta magnitud. Además, también se debe establecer cómo
varía el tamaño de los cursos a medida que se avanza en las carreras, ya
que es esperable que este vaya disminuyendo. En la
Figura~\ref{fig-semestres} se muestra una caída pronunciada y sostenida
en el número de estudiantes por curso a medida que avanza el semestre
curricular. Los primeros cuatro semestres mantienen tamaños
relativamente grandes y estables (\textasciitilde50--54 estudiantes), lo
que refleja la estructura de cursos obligatorios masivos propios de los
ciclos iniciales de las carreras. A partir del quinto semestre se
produce un quiebre, con una caída que lleva el promedio de
\textasciitilde53 a \textasciitilde35 estudiantes, y que continúa de
forma acelerada hasta estabilizarse en valores muy bajos en los
semestres finales. Esta trayectoria es consistente con la lógica
curricular esperada: los semestres avanzados concentran cursos
electivos, seminarios y actividades de titulación de menor matrícula.

\begin{figure}

\centering{

\pandocbounded{\includegraphics[keepaspectratio]{04-resultados-curso_files/figure-pdf/fig-tamano-1.pdf}}

}

\caption{\label{fig-tamano}Distribución de notas de cursos pequeños y
normales FACSO (2018-2025)}

\end{figure}%

Con el fin de despejar incógnitas, la Figura~\ref{fig-tamano} muestra la
distribución de notas para cursos con menos de 10 estudiantes y para
cursos con 10 o más estudiantes. Ambos gráficos muestran una
concentración de notas en los rangos altos de la escala, pero las
diferencias entre ambas distribuciones son sustantivas. Los cursos con
menos de 10 estudiantes presentan una distribución fuertemente
concentrada en las notas máximas ---especialmente en torno a 7,0--- y
una marcada asimetría, junto con algunos casos extremos de notas bajas,
lo que sugiere una menor variabilidad interna y una tendencia a
calificaciones más altas. En cambio, los cursos con 10 o más estudiantes
exhiben una distribución más equilibrada y dispersa, con una forma
cercana a la normal y una concentración principal entre 5,5 y 6,5. Así,
mientras los cursos pequeños tienden a acumular calificaciones
excepcionalmente altas, los cursos más numerosos muestran patrones de
evaluación más heterogéneos y moderados.

Estos cursos pequeños, en la mayoría de los casos, son cursos con sólo
un estudiante, puesto que son ramos donde las personas desarrollan sus
prácticas profesionales y/o memorias de título. El problema con esto es
que un curso de una persona tendría el mismo peso que un curso de 100
estudiantes, por ejemplo, lo cual distorsionaría en gran medida los
resultados de los análisis. Con el objetivo de realizar un estudio con
el menor margen de error, en este informe se toma la decisión de dejar
fuera de los análisis posteriores todo curso que tenga menos de 10
estudiantes, lo que corresponde a 2368 cursos.

\subsection{Notas}\label{notas}

\subsubsection{Nivel facultad}\label{nivel-facultad}

\begin{figure}

\centering{

\pandocbounded{\includegraphics[keepaspectratio]{04-resultados-curso_files/figure-pdf/fig-facso-notas-1.pdf}}

}

\caption{\label{fig-facso-notas}Elaboración propia (2026)}

\end{figure}%

La Figura~\ref{fig-facso-notas} presenta la distribución de las notas a
nivel facultad una vez eliminados los cursos con menos de 10
estudiantes. El gráfico muestra una distribución de notas con un
promedio 6.0, difiriendo con el resultado antes del filtrado. Su
distribución se concentra principalmente entre 5,5 y 6,5, representando
un 67\% del total, y con su punto máximo en torno a 6,1--6,3, lo que
indica un predominio de cursos con calificaciones altas en la facultad.
La distribución presenta una leve asimetría hacia la izquierda,
evidenciada por una cola de casos con notas más bajas, aunque estos son
relativamente poco frecuentes. Asimismo, las notas extremas bajo 5,0
aparecen de manera marginal, mientras que las calificaciones cercanas a
7,0 disminuyen progresivamente. En conjunto, los resultados sugieren
calificaciones elevadas y homogéneas.

\subsubsection{Facultad y carreras a lo largo del
tiempo}\label{facultad-y-carreras-a-lo-largo-del-tiempo}

\begin{figure}

\centering{

\pandocbounded{\includegraphics[keepaspectratio]{04-resultados-curso_files/figure-pdf/fig-facso-notas-tiempo-1.pdf}}

}

\caption{\label{fig-facso-notas-tiempo}Elaboración propia (2026)}

\end{figure}%

En Figura~\ref{fig-facso-notas-tiempo} se presentan las notas promedio
por año a nivel facultad junto a otras medidas estadísticas. El gráfico
muestra un aumento sostenido de las notas promedio de la facultad entre
2018 y 2021, pasando desde medianas cercanas a 5,5--5,7 hacia valores
sobre 6,1, lo que sugiere un incremento generalizado en las
calificaciones, bien la flexibilización de los estándares de evaluación
durante ese período. A partir de 2021, las distribuciones tienden a
estabilizarse en torno a promedios altos, manteniéndose relativamente
constantes hasta 2025, aunque con una leve disminución en 2024 y una
recuperación parcial en 2025. Asimismo, desde 2020 en adelante se
observa una mayor dispersión y una presencia más frecuente de valores
atípicos bajos, especialmente entre 2022 y 2025, indicando que, pese a
la estabilidad de las medianas, persisten diferencias importantes en el
desempeño de algunos cursos o estudiantes.

\begin{figure}

\centering{

\pandocbounded{\includegraphics[keepaspectratio]{04-resultados-curso_files/figure-pdf/fig-box-notas-facso-1.pdf}}

}

\caption{\label{fig-box-notas-facso}Elaboración propia (2026)}

\end{figure}%

En Figura~\ref{fig-box-notas-facso} se pueden observar los promedios de
las distintas carreras considerando información entre 2018 y 2025. El
gráfico evidencia diferencias moderadas pero consistentes entre carreras
en la distribución de las notas. Psicología y Educación Parvularia
presentan las medianas más altas, ambas sobre 6,2, mientras
Antropología, Sociología y Trabajo Social muestran medianas alrededor de
6,0. Sociología y Antropología destacan además por exhibir una mayor
dispersión y una mayor cantidad de valores atípicos bajos, indicando
heterogeneidad en sus cursos. Por su parte, Educación Parvularia
presenta una distribución más compacta y homogénea, mientras que
Psicología combina medianas altas con una dispersión relativamente
amplia. En conjunto, aunque todas las carreras concentran sus notas en
niveles elevados, existen diferencias relevantes en términos de
variabilidad de los desempeños académicos.

\begin{figure}

\centering{

\pandocbounded{\includegraphics[keepaspectratio]{04-resultados-curso_files/figure-pdf/fig-trayec-notas-facso-1.pdf}}

}

\caption{\label{fig-trayec-notas-facso}Elaboración propia (2026)}

\end{figure}%

Figura~\ref{fig-trayec-notas-facso} presenta la evolución de las notas
promedio por carrera entre 2018 y 2025. En términos generales, todas las
carreras muestran un patrón similar: las notas aumentan entre 2018 y
2020-2021, para luego estabilizarse o descender levemente hacia el final
del período. Este comportamiento es especialmente claro en Psicología y
Educación Parvularia, que mantienen de manera consistente los promedios
más altos durante todo el período, superando en la mayoría de los años
el 6,0. En contraste, Antropología presenta una trayectoria de menor
rendimiento, partiendo desde 5,6 en 2018, alcanzando un peak en 2020
cercano al 6,0, y mostrando un descenso sostenido hacia 2025 hasta
situarse en torno al 5,7. Sociología y Trabajo Social exhiben
trayectorias intermedias y relativamente similares entre sí, con valores
que oscilan entre 5,5 y 6,0 a lo largo del período, aunque con algunas
fluctuaciones puntuales como la caída de Sociología en 2019.

Un elemento destacable es el alza generalizada observada entre 2019 y
2021, presente en todas las carreras, lo que podría asociarse al
contexto de pandemia y los ajustes en las modalidades de evaluación. En
conjunto, los resultados evidencian una convergencia general hacia
niveles altos de notas, con Educación Parvularia y Psicología liderando
consistentemente el período, y Antropología como la carrera con el
rendimiento promedio más bajo y la tendencia más descendente en los
últimos años.

\subsection{Tasa de aprobación}\label{tasa-de-aprobaciuxf3n}

Otra arista importante para entender las calificaciones en la facultad,
y que está directamente ligada a las notas, es el porcentaje de
aprobación dentro de los cursos. El porcentaje de aprobación puede
entregar pistas de la complejidad que implica cursar ramos
universitarios y superarlos exitosamente. En este sentido, el aprobar un
ramo debería ser signo de reconocimiento respecto al cumplimiento de los
aprendizajes esperados dentro de una cátedra.

\subsubsection{Facultad en el tiempo}\label{facultad-en-el-tiempo}

\begin{longtable}[t]{lrr}

\caption{\label{tbl-aprobacion}Porcentaje de aprobación de la facultad
por año:}

\tabularnewline

\\
\toprule
\cellcolor[HTML]{952626}{\textcolor{white}{\textbf{Año}}} & \cellcolor[HTML]{952626}{\textcolor{white}{\textbf{Aprobación promedio (\%)}}} & \cellcolor[HTML]{952626}{\textcolor{white}{\textbf{N cursos}}}\\
\midrule
2018 & 97.19 & 77\\
2019 & 97.50 & 145\\
2020 & 97.80 & 273\\
2021 & 97.49 & 369\\
2022 & 97.11 & 381\\
\addlinespace
2023 & 97.15 & 392\\
2024 & 97.01 & 367\\
2025 & 97.21 & 287\\
\cellcolor[HTML]{f2f2f2}{\textbf{Total}} & \cellcolor[HTML]{f2f2f2}{\textbf{97.31}} & \cellcolor[HTML]{f2f2f2}{\textbf{2291}}\\
\bottomrule

\end{longtable}

La Tabla~\ref{tbl-aprobacion} presenta los porcentajes de aprobación en
FACSO por año. Se puede observar un patrón estable en el tiempo: todos
los años tienen un 97\% de aprobación promedio, variando solamente en
algunas décimas. A pesar de que 2022 y 2023 concentran mayores
cantidades de cursos que otros años, la consistencia en la aprobación de
los cursos es persistente. Este resultado indica que en FACSO existe una
cantidad mínima de cursos que cuentan con estudiantes reprobados, lo
cual también se puede comprender como una alta homogeneización en
términos evaluativos al notar que solamente un 3\% de la población
estudiantil reprueba.

\subsubsection{Carreras en el tiempo}\label{carreras-en-el-tiempo}

\begin{figure}

\centering{

\pandocbounded{\includegraphics[keepaspectratio]{04-resultados-curso_files/figure-pdf/fig-aprobacion-1.pdf}}

}

\caption{\label{fig-aprobacion}Elaboración propia (2026)}

\end{figure}%

En la Figura~\ref{fig-aprobacion} se muestran los porcentajes de
aprobación promedio de cada carrera desde el 2018 hasta el 2025. El
gráfico muestra que las tasas de aprobación en todas las carreras de
FACSO se mantienen consistentemente elevadas a lo largo del período
2018--2025, oscilando casi siempre entre el 95\% y el 100\%, lo que
indica que la reprobación es un fenómeno marginal en la facultad con
independencia de la carrera o el año. Dicho esto, emergen algunas
diferencias relevantes: Psicología y Educación Parvularia son las
carreras más estables y con tasas sostenidamente altas, rozando el 99\%
en varios años. Antropología, en cambio, muestra la trayectoria más
volátil y una tendencia levemente decreciente, partiendo cerca del 100\%
en 2018 y situándose en torno al 96\% hacia 2024--2025. Por su parte,
Sociología es la carrera con mayor variabilidad entre años y registra
los valores más bajos del conjunto, con una caída en 2019
(\textasciitilde94\%) y fluctuaciones persistentes que la ubican por
debajo del resto de carreras. Trabajo Social presenta un comportamiento
intermedio, con ciertas variaciones pero relativamente estable. En suma,
aunque el nivel general de aprobación es alto en toda la facultad,
Sociología y Antropología exhiben una fragilidad relativa que las
distingue del resto al ser durante varios años las carreras con las
tasas de aprobación más bajas, aunque el panorama general sigue
mostrando una alta aprobación en todas las carreras.

\subsection{Dispersión de notas}\label{dispersiuxf3n-de-notas}

Hasta ahora, sabemos que las notas individuales son altas, y que la tasa
de reprobación es muy baja, pero no sabemos qué tanto varían las notas
dentro de cada curso. Esto es especialmente importante, ya que cursos
con una dispersión más homogénea no permiten distinguir a los y las
estudiantes que tienen un rendimiento sobresaliente de aquellos con un
rendimiento regular, inhibiendo el rol diferenciador de las
evaluaciones, lo que puede causar una sensación de injusticia en la
población estudiantil. Por ello, además de estudiar las calificaciones,
es relevante entender la dispersión de las notas.

La desviación estándar es una medida estadística que permite evaluar el
grado de dispersión o variabilidad de las notas respecto de su promedio.
En términos simples, indica qué tan homogéneo o heterogéneo es el
desempeño académico de los y las estudiantes dentro de un curso. Por
ejemplo, si el promedio de notas de la facultad es 5.9 y tiene una
desviación estándar de 1, esto significa el 68\% de los casos se
encuentra entre +1 y -1 desviación estándar, es decir, entre el 5.8 y el
6.0 (asumiendo una distribución normal de las notas).

El análisis de esta medida resulta especialmente relevante para
comprender las dinámicas de evaluación al interior de los cursos, ya que
permite complementar la información entregada por los promedios de
notas. De esta manera, la desviación estándar contribuye a identificar
cursos con evaluaciones más homogéneas o más diferenciadoras, entregando
una visión más completa sobre las prácticas de evaluación y el desempeño
estudiantil.

\subsubsection{Nivel facultad}\label{nivel-facultad-1}

\begin{figure}

\centering{

\pandocbounded{\includegraphics[keepaspectratio]{04-resultados-curso_files/figure-pdf/fig-desv-facso-1.pdf}}

}

\caption{\label{fig-desv-facso}Elaboración propia (2026)}

\end{figure}%

En la Figura~\ref{fig-desv-facso} se muestra la concentración de las
desviaciones estándar de los cursos impartidos en FACSO. Se puede
observar que la mayoría de los cursos presentan desviaciones estándar
entre 0 y 1, donde la mayor concentración se ubica en el 0.4. Esto
quiere decir que el patrón más recurrente es que los cursos tienden a
variar + o - 0.4 décimas respecto a su promedio. Con ello, los
resultados indican que la mayoría de los cursos en FACSO mantienen una
variabilidad media baja.

Considerando que el promedio de notas de la facultad es 6.0, una
desviación estándar de 0.4 implica que aproximadamente el 68\% de los
estudiantes en un curso típico se encuentra entre 5.6 y 6.4, mientras
que el 96\% se halla entre 5.2 y 6.8, lo que sugiere una alta
homogeneidad en las calificaciones. Sin embargo, también se observa una
cola hacia desviaciones estándar más altas, lo que indica que existen
cursos con mayor heterogeneidad en las notas, aunque estos son menos
frecuentes.

\subsubsection{Facultad en el tiempo}\label{facultad-en-el-tiempo-1}

\begin{figure}

\centering{

\pandocbounded{\includegraphics[keepaspectratio]{04-resultados-curso_files/figure-pdf/fig-desv-facso-tiempo-1.pdf}}

}

\caption{\label{fig-desv-facso-tiempo}Elaboración propia (2026)}

\end{figure}%

En la Figura~\ref{fig-desv-facso-tiempo} se presentan las desviaciones
estándar promedios por año a nivel facultad. En el gráfico se puede
observar que el 2019 es el año con mayor dispersión de notas, seguido de
una baja el 2020 que posteriormente vuelve a tomar un patrón ascendente
hasta 2022. Seguido de esto, la trayectoria presenta un comportamiento
irregular donde 2023 disminuye la dispersión, el año siguiente aumenta
y, finalmente, el 2025 se caracteriza por ser el año con mayor
homogeneidad en las calificaciones. Con todo, se puede evidenciar que a
medida que pasa el tiempo las evaluaciones tienden a homogeneizarse, lo
cual se expresa en una menor dispersión promedio.

\subsubsection{Nivel carrera}\label{nivel-carrera}

\begin{figure}

\centering{

\pandocbounded{\includegraphics[keepaspectratio]{04-resultados-curso_files/figure-pdf/fig-desv-carr-1.pdf}}

}

\caption{\label{fig-desv-carr}Elaboración propia (2026)}

\end{figure}%

En la Figura~\ref{fig-desv-carr} se muestran las concentraciones de
desviaciones estándar por carrera. En el gráfico se puede observar que
todas las carreras comparten un patrón similar: la gran mayoría de sus
cursos concentra desviaciones estándar bajas (entre 0.2 y 0.5), lo que
indica que en general el rendimiento estudiantil tiende a ser homogéneo.
Sin embargo, existen diferencias relevantes en las colas de la
distribución: Antropología destaca claramente del resto al presentar una
cola derecha notoriamente más extensa, con cursos que alcanzan
desviaciones estándar de hasta 2.0--2.5, evidenciando que poseen una
proporción mayor de cursos con rendimiento altamente heterogéneo.
Educación Parvularia y Trabajo Social, en cambio, muestran las
distribuciones más concentradas y con colas más cortas, sugiriendo mayor
uniformidad en sus calificaciones. Psicología presenta la forma más alta
y estrecha, lo que la posiciona como la carrera con menor dispersión
interna en sus calificaciones. Por su parte, Sociología ocupa una
posición intermedia con una dispersión moderada, pero con una cola
derecha más amplia que las demás carreras, sugiriendo que presenta una
mayor cantidad de cursos con desviaciones estándar mayores a 1. En
síntesis, si bien el patrón general apunta hacia notas homogéneas en
toda la facultad, Antropología y Sociología resaltan como la excepción,
con una variabilidad interna mayor que el resto de las carreras.

\subsubsection{Carreras en el tiempo}\label{carreras-en-el-tiempo-1}

\begin{figure}

\centering{

\pandocbounded{\includegraphics[keepaspectratio]{04-resultados-curso_files/figure-pdf/fig-desv-carr-tiempo-1.pdf}}

}

\caption{\label{fig-desv-carr-tiempo}Elaboración propia (2026)}

\end{figure}%

En la Figura~\ref{fig-desv-carr-tiempo} se muestra la evolución de la
dispersión de notas por carrera entre 2018 y 2025, revelando
trayectorias claramente diferenciadas. Psicología destaca por una
tendencia decreciente sostenida, convergiendo hacia valores cercanos a
0.43--0.47 en 2025, lo que indica una progresiva homogeneización del
rendimiento. Si bien Educación Parvularia también presenta una
trayectoria descendente, tiene la particularidad que el año 2024 su
dispersión se dispara en términos heterogéneos para luego decaer
abruptamente el año siguiente. Sociología, en contraste, es la carrera
con mayor y más persistente dispersión a lo largo de todo el período,
alcanzando su punto más alto el 2019 y niveles que se mantienen elevados
hasta 2022 para luego descender moderadamente, confirmando la
heterogeneidad en los cursos que la carrera. Por su lado, Antropología
se presenta con una tendencia ascendente entre 2018 y 2024 que los años
posteriores decae, mientras que Trabajo Social tiene sus puntos de
dispersión más altos los años límites del estudio, notándose un patrón
homogeneización entre 2020 y 2023. En conjunto, los datos sugieren que
la dispersión en FACSO no es un fenómeno uniforme sino que responde a
dinámicas propias de cada carrera.

\section{Cruces}\label{cruces}

Esta subsección tiene el objetivo de profundizar en los resultados ya
encontrados a partir de la asociación de las calificaciones de los
cursos con variables que caracterizan a estos mismos. En esta línea, el
paso siguiente es comprender en qué medida el tipo de curso, el género y
el tamaño del curso determinan las notas universitarias.

\subsection{Tipo de curso}\label{tipo-de-curso}

\subsubsection{Nivel facultad}\label{nivel-facultad-2}

\begin{figure}

\centering{

\pandocbounded{\includegraphics[keepaspectratio]{04-resultados-curso_files/figure-pdf/fig-car-facso-1.pdf}}

}

\caption{\label{fig-car-facso}Elaboración propia (2026)}

\end{figure}%

En la Figura~\ref{fig-car-facso} se presentan las distribuciones de
notas distinguiendo por cursos obligatorios y electivos. El gráfico
muestra que los cursos electivos tienen una mediana 6.3 con una
distribución compacta y con algunos casos extremos. Por su parte, los
cursos obligatorios poseen una mediana cercana al 6.0 con una
distribución ligeramente más dispersa que los cursos electivos, pero con
una mayor cantidad de cursos que presentan valores muy bajos. En
síntesis, los cursos electivos tienen valores más altos y concentrados
en la parte alta de la escala de notas en comparación a los cursos
obligatorios.

\subsubsection{Trayectoria nivel
facultad}\label{trayectoria-nivel-facultad}

\begin{figure}

\centering{

\pandocbounded{\includegraphics[keepaspectratio]{04-resultados-curso_files/figure-pdf/fig-car-facso-tiempo-1.pdf}}

}

\caption{\label{fig-car-facso-tiempo}Elaboración propia (2026)}

\end{figure}%

En la Figura~\ref{fig-car-facso-tiempo} se muestran los promedios de
notas de los cursos electivos obligatorios por año. Los resultados
muestran que durante todo el período analizado los cursos electivos
mantienen notas promedio consistentemente superiores a los obligatorios.
Ambas categorías experimentan un alza notable en 2020--2021, atribuible
a la flexibilidad académica durante la pandemia, pero la trayectoria
post-pandemia revela una divergencia estructural: mientras los electivos
se estabilizan en torno a 6.2--6.3, los cursos obligatorios muestran una
caída sostenida que hacia 2025 los sitúa cerca a los niveles
pre-pandemia (\textasciitilde5.8), sugiriendo que el retorno a la
presencialidad impactó de forma asimétrica ambos tipos de cursos.

\subsubsection{Nivel carrera}\label{nivel-carrera-1}

\subsubsection{Sociología}

\begin{figure}

\centering{

\pandocbounded{\includegraphics[keepaspectratio]{04-resultados-curso_files/figure-pdf/fig-box-socio-1.pdf}}

}

\caption{\label{fig-box-socio}Elaboración propia (2026)}

\end{figure}%

\subsubsection{Psicología}

\begin{figure}

\centering{

\pandocbounded{\includegraphics[keepaspectratio]{04-resultados-curso_files/figure-pdf/fig-box-psico-1.pdf}}

}

\caption{\label{fig-box-psico}Elaboración propia (2026)}

\end{figure}%

\subsubsection{Antropología}

\begin{figure}

\centering{

\pandocbounded{\includegraphics[keepaspectratio]{04-resultados-curso_files/figure-pdf/fig-box-ant-1.pdf}}

}

\caption{\label{fig-box-ant}Elaboración propia (2026)}

\end{figure}%

\subsubsection{Trabajo Social}

\begin{figure}

\centering{

\pandocbounded{\includegraphics[keepaspectratio]{04-resultados-curso_files/figure-pdf/fig-box-ts-1.pdf}}

}

\caption{\label{fig-box-ts}Elaboración propia (2026)}

\end{figure}%

\subsubsection{Educación Parvularia}

\begin{figure}

\centering{

\pandocbounded{\includegraphics[keepaspectratio]{04-resultados-curso_files/figure-pdf/fig-box-ed-1.pdf}}

}

\caption{\label{fig-box-ed}Elaboración propia (2026)}

\end{figure}%

En este panel se observan las distribuciones de las notas distinguiendo
cursos electivos y obligatorios por carrera. Los boxplots muestran un
patrón consistente en las 5 carreras de FACSO: los cursos electivos
presentan medianas levemente superiores exceptuando Educación
Parvularia, con cajas estrechas y prácticamente sin valores atípicos en
comparación a los cursos obligatorios, lo que refleja calificaciones
homogéneas concentradas en el tramo alto. Los cursos obligatorios, en
cambio, exhiben mayor variabilidad y outliers hacia la baja en todas las
unidades académicas ---particularmente en Antropología y Educación
Parvularia---, evidenciando una distribución más heterogénea con
presencia de rendimientos inferiores. Con todo, el hecho de que este
patrón se replique de forma casi idéntica en Sociología, Psicología,
Antropología y Trabajo Social sugiere que se trata de un fenómeno
estructural de la facultad y no de una particularidad de alguna carrera
específica.

\subsubsection{Trayectorias por carrera}\label{trayectorias-por-carrera}

\subsubsection{Sociología}

\begin{figure}

\centering{

\pandocbounded{\includegraphics[keepaspectratio]{04-resultados-curso_files/figure-pdf/fig-car-socio-t-1.pdf}}

}

\caption{\label{fig-car-socio-t}Elaboración propia (2026)}

\end{figure}%

\subsubsection{Psicología}

\begin{figure}

\centering{

\pandocbounded{\includegraphics[keepaspectratio]{04-resultados-curso_files/figure-pdf/fig-car-psico-t-1.pdf}}

}

\caption{\label{fig-car-psico-t}Elaboración propia (2026)}

\end{figure}%

\subsubsection{Antropología}

\begin{figure}

\centering{

\pandocbounded{\includegraphics[keepaspectratio]{04-resultados-curso_files/figure-pdf/fig-car-ant-t-1.pdf}}

}

\caption{\label{fig-car-ant-t}Elaboración propia (2026)}

\end{figure}%

\subsubsection{Trabajo Social}

\begin{figure}

\centering{

\pandocbounded{\includegraphics[keepaspectratio]{04-resultados-curso_files/figure-pdf/fig-car-ts-t-1.pdf}}

}

\caption{\label{fig-car-ts-t}Elaboración propia (2026)}

\end{figure}%

\subsubsection{Educación Parvularia}

\begin{figure}

\centering{

\pandocbounded{\includegraphics[keepaspectratio]{04-resultados-curso_files/figure-pdf/fig-car-ed-t-1.pdf}}

}

\caption{\label{fig-car-ed-t}Elaboración propia (2026)}

\end{figure}%

Este panel de gráficos muestra la trayectoria de notas promedio por
carácter del curso en cada carrera. Los electivos superan
consistentemente a los obligatorios en casi todas las carreras, pero con
matices importantes. Sociología mantiene una brecha estable y
persistente a lo largo de todo el período, con los obligatorios
mostrando una caída hacia 2025 que amplía la diferencia respecto a los
electivos. Asimismo, Psicología presenta la brecha más pronunciada de la
facultad, con electivos que se estabilizan en torno a 6.4 mientras los
obligatorios descienden sostenidamente desde 2020, llegando a
\textasciitilde5.8 en 2025, lo que configura una divergencia
post-pandemia particularmente marcada. Antropología es el caso más
llamativo: ambas líneas convergen casi completamente desde 2020 en
adelante, rondando el 5.8--6.0, sugiriendo que en esta carrera el
carácter del curso no genera diferencias relevantes en el rendimiento.
Por su parte, Trabajo Social muestra una brecha moderada pero con
obligatorios que se estabilizan cerca del 5.9, mientras los electivos
tienen un patrón estable sobre el 6.0. Educación Parvularia vuelve a
destacar como caso atípico: los obligatorios parten más altos que los
electivos en 2018 y ambas líneas se mantienen entrelazadas durante todo
el período, con los electivos superando levemente el 2025
(\textasciitilde6.6 vs \textasciitilde6.2), reforzando la idea de que en
esta carrera el carácter del curso tiene escasa incidencia en las
calificaciones.

\subsection{Género}\label{guxe9nero-1}

\subsubsection{Nivel facultad}\label{nivel-facultad-3}

\begin{figure}

\centering{

\pandocbounded{\includegraphics[keepaspectratio]{04-resultados-curso_files/figure-pdf/fig-gen-facso-1.pdf}}

}

\caption{\label{fig-gen-facso}Elaboración propia (2026)}

\end{figure}%

La Figura~\ref{fig-gen-facso} muestra una asociación leve pero visible
entre la composición de género y las notas: los cursos con mayor
proporción de mujeres (81--100\%) presentan las medianas más altas
(\textasciitilde6.3) y menor dispersión, mientras que los grupos de
composición mixta (21--80\%) exhiben medianas más bajas
(\textasciitilde5.9--6.0) y mayor variabilidad, con outliers que
descienden hasta notas cercanas a 4.0. Aun así, las diferencias
absolutas son pequeñas y deben interpretarse con cautela.

\subsubsection{Facultad en el tiempo}\label{facultad-en-el-tiempo-2}

\begin{figure}

\centering{

\pandocbounded{\includegraphics[keepaspectratio]{04-resultados-curso_files/figure-pdf/fig-gen-facso-tiempo-1.pdf}}

}

\caption{\label{fig-gen-facso-tiempo}Elaboración propia (2026)}

\end{figure}%

En la Figura~\ref{fig-gen-facso-tiempo} se presentan las trayectorias de
los cursos agrupando por composición de género. El gráfico confirma y
enriquece el hallazgo anterior: los cursos con mayor proporción de
mujeres (81--100\%) mantienen una ventaja sostenida y estable a lo largo
de todo el período 2018--2025, situándose consistentemente por encima
del resto (a excepción del 2024) con notas que oscilan entre 5.8 y 6.4.
El resto de los grupos muestra trayectorias más entrelazadas y
volátiles, con un alza generalizada en 2020--2021 que podría estar
relacionada al efecto pandemia y una posterior convergencia hacia
valores similares (\textasciitilde5.8--6.0) desde 2022 en adelante. Se
destaca que los grupos de composición mixta (41--60\% y 21--40\%) no
solo parten más bajos sino que exhiben mayor inestabilidad en sus
trayectorias, mientras que el grupo 61--80\% se comporta de forma más
estable pero sin alcanzar las calificaciones promedio del grupo de mayor
feminización. En síntesis, la composición de género predominantemente
femenina se asocia con calificaciones superiores y más estables en el
tiempo, mientras que los cursos con una mayor heterogeneidad de género
presentan notas más bajas y trayectorias más cambiantes.

\subsubsection{Nivel carrera}\label{nivel-carrera-2}

\subsubsection{Sociología}

\begin{figure}

\centering{

\pandocbounded{\includegraphics[keepaspectratio]{04-resultados-curso_files/figure-pdf/fig-gen-socio-1.pdf}}

}

\caption{\label{fig-gen-socio}Elaboración propia (2026)}

\end{figure}%

\subsubsection{Psicología}

\begin{figure}

\centering{

\pandocbounded{\includegraphics[keepaspectratio]{04-resultados-curso_files/figure-pdf/fig-gen-psico-1.pdf}}

}

\caption{\label{fig-gen-psico}Elaboración propia (2026)}

\end{figure}%

\subsubsection{Antropología}

\begin{figure}

\centering{

\pandocbounded{\includegraphics[keepaspectratio]{04-resultados-curso_files/figure-pdf/fig-gen-ant-1.pdf}}

}

\caption{\label{fig-gen-ant}Elaboración propia (2026)}

\end{figure}%

\subsubsection{Trabajo Social}

\begin{figure}

\centering{

\pandocbounded{\includegraphics[keepaspectratio]{04-resultados-curso_files/figure-pdf/fig-gen-ts-1.pdf}}

}

\caption{\label{fig-gen-ts}Elaboración propia (2026)}

\end{figure}%

\subsubsection{Educación Parvularia}

\begin{figure}

\centering{

\pandocbounded{\includegraphics[keepaspectratio]{04-resultados-curso_files/figure-pdf/fig-gen-ed-1.pdf}}

}

\caption{\label{fig-gen-ed}Elaboración propia (2026)}

\end{figure}%

Los gráficos de este panel desagregan la relación entre composición de
género de los cursos y notas por carrera. Al observar una por una las
figuras, se puede notar la falta de categorías para algunas carreras
como Trabajo Social y Educación Parvularia, lo cual se debe a la alta
matrícula femenina en detrimento de la masculina en estos programas,
limitando la comparación con las demás carreras. En Sociología se
aprecia una leve ventaja del grupo 81--100\% (\textasciitilde6.4)
respecto a los demás, aunque con dispersión considerable y
calificaciones inferiores en los grupos intermedios. En Psicología las
distribuciones son prácticamente idénticas entre todos los grupos pero
difieren en la mediana; se aprecia que el grupo de menor presencia
femenina es apenas visible, lo que puede deberse a la poca cantidad de
cursos dentro de esta categoría. En Antropología todas las categorías
disponibles muestran dispersiones similares, pero es llamativo que la
categoría más femenizada es la que tiene la mediana más baja. Por su
lado, Trabajo Social, que solo cuenta con tres grupos, presenta un
patrón escalonado ascendente, donde los cursos con mayor presencia
femenina tienen mejores calificaciones. Educación Parvularia presenta
únicamente el grupo 81--100\%, lo que impide cualquier comparación
interna. Si bien se pueden observar diferencias en los patrones entre
carreras, el principal hallazgo es la inexistencia de cursos con
proporciones bajas de mujeres en carreras como Trabajo Social y
Educación Parvularia sugiriendo una predominación femenina de
estudiantes que deciden cursar estos programas.

\subsubsection{Carrera en el tiempo}\label{carrera-en-el-tiempo}

\subsubsection{Sociología}

\begin{figure}

\centering{

\pandocbounded{\includegraphics[keepaspectratio]{04-resultados-curso_files/figure-pdf/fig-gen-socio-t-1.pdf}}

}

\caption{\label{fig-gen-socio-t}Elaboración propia (2026)}

\end{figure}%

\subsubsection{Psicología}

\begin{figure}

\centering{

\pandocbounded{\includegraphics[keepaspectratio]{04-resultados-curso_files/figure-pdf/fig-gen-psico-t-1.pdf}}

}

\caption{\label{fig-gen-psico-t}Elaboración propia (2026)}

\end{figure}%

\subsubsection{Antropología}

\begin{figure}

\centering{

\pandocbounded{\includegraphics[keepaspectratio]{04-resultados-curso_files/figure-pdf/fig-gen-ant-t-1.pdf}}

}

\caption{\label{fig-gen-ant-t}Elaboración propia (2026)}

\end{figure}%

\subsubsection{Trabajo Social}

\begin{figure}

\centering{

\pandocbounded{\includegraphics[keepaspectratio]{04-resultados-curso_files/figure-pdf/fig-gen-ts-t-1.pdf}}

}

\caption{\label{fig-gen-ts-t}Elaboración propia (2026)}

\end{figure}%

\subsubsection{Educación Parvularia}

\begin{figure}

\centering{

\pandocbounded{\includegraphics[keepaspectratio]{04-resultados-curso_files/figure-pdf/fig-gen-ed-t-1.pdf}}

}

\caption{\label{fig-gen-ed-t}Elaboración propia (2026)}

\end{figure}%

El anterior grupo de gráficos muestra la trayectoria de notas según
composición de género al interior de cada carrera, y el hallazgo
transversal más relevante es que no existe un patrón consistente ni
uniforme entre carreras. En Sociología, el grupo 81--100\% mujeres
destaca con notas más altas y estables (\textasciitilde6.4--6.6),
mientras el resto de los grupos presenta un alza notable en 2020 para
luego variar con mayor volatilidad en torno al 6.0. En Psicología, las
líneas se entrelazan de forma irregular y el grupo 21--40\% muestra una
caída abrupta en 2023 (\textasciitilde5.0) que distorsiona cualquier
tendencia, sin que ningún grupo muestre una ventaja sostenida clara. Por
su parte, en Antropología las categorías intermedias oscilan en un rango
estrecho (\textasciitilde5.8--6.1) sin diferencias sistemáticas, pero se
observa un resultado llamativo: desde 2024 se invierte el patrón y los
cursos con menor presencia femenina son aquellos con mejores
calificaciones. En Trabajo Social, el grupo 81--100\% muestra una
ventaja más visible y creciente hacia 2025 (\textasciitilde6.3),
mientras el grupo 41--60\% cae notoriamente en los últimos años.
Educación Parvularia, con un único grupo disponible, se mantiene estable
en torno a 6.1--6.3. En síntesis, si bien hay una tendencia en que la
composición femenina obtiene mejores calificaciones, esto no se explica
de forma sistemática, ya que las variaciones observadas responden más a
dinámicas propias de cada carrera que a un efecto atribuible al género.

\subsection{Cantidad de estudiantes}\label{cantidad-de-estudiantes}

\subsubsection{Nivel facultad}\label{nivel-facultad-4}

\begin{figure}

\centering{

\pandocbounded{\includegraphics[keepaspectratio]{04-resultados-curso_files/figure-pdf/fig-est-facso-1.pdf}}

}

\caption{\label{fig-est-facso}Elaboración propia (2026)}

\end{figure}%

En la Figura~\ref{fig-est-facso} se observa una relación inversa leve
entre el tamaño del curso y el rendimiento: los cursos más pequeños (≤20
y 21--40 estudiantes) presentan medianas ligeramente más altas
(\textasciitilde6.1--6.2) y cajas más compactas en el rango superior,
mientras que los cursos medianos y grandes (41--60, 61--80 y ≥81)
muestran medianas algo menores (\textasciitilde5.9--6.0) y mayor
dispersión, con cajas que se extienden hacia valores más bajos. Sin
embargo, las diferencias absolutas entre grupos son pequeñas y todos los
grupos comparten outliers hacia la baja, siendo los cursos de tamaño
medio (41--60) los que presentan los valores atípicos más extremos
(\textasciitilde3.9). En síntesis, aunque existe una tendencia a que los
cursos más pequeños obtengan notas ligeramente superiores y más
homogéneas, las diferencias no llegan a ser concluyentes.

\subsubsection{Facultad por año}\label{facultad-por-auxf1o}

\begin{figure}

\centering{

\pandocbounded{\includegraphics[keepaspectratio]{04-resultados-curso_files/figure-pdf/fig-est-facso-tiempo-1.pdf}}

}

\caption{\label{fig-est-facso-tiempo}Elaboración propia (2026)}

\end{figure}%

La Figura~\ref{fig-est-facso-tiempo} muestra la trayectoria de notas
según tamaño del curso entre 2018 y 2025, y el hallazgo más relevante es
que las diferencias inicialmente observadas entre grupos se diluyen con
el tiempo. En 2018 existe una diferenciación clara, con los cursos más
pequeños (≤20) partiendo considerablemente más alto (\textasciitilde6.3)
mientras los cursos grandes (61--80) tienen notas más bajas
(\textasciitilde5.5), pero desde 2020 todas las líneas convergen en un
rango estrecho de 5.9--6.2, sugiriendo que el efecto pandemia
homogeneizó el rendimiento independientemente del tamaño. Desde el 2023
la dinámica se estabiliza, donde todos los grupos muestran una tendencia
levemente decreciente (a excepción de los cursos entre 61-80
estudiantes) que los sitúa en torno a 5.8--6.2 hacia 2025, pero los
cursos pequeños (≤20 y 21--40) se mantienen ligeramente con mejores
calificaciones. En síntesis, el tamaño del curso muestra una asociación
débil y decreciente con las notas: la ventaja de los cursos pequeños,
notoria en 2018, se ha ido erosionando progresivamente, y la trayectoria
post-pandemia apunta hacia una convergencia general con leve desventaja
de los cursos más grandes.

\subsubsection{Nivel carrera}\label{nivel-carrera-3}

\subsubsection{Sociología}

\begin{figure}

\centering{

\pandocbounded{\includegraphics[keepaspectratio]{04-resultados-curso_files/figure-pdf/fig-n-carrera-socio-1.pdf}}

}

\caption{\label{fig-n-carrera-socio}Elaboración propia (2026)}

\end{figure}%

\subsubsection{Psicología}

\begin{figure}

\centering{

\pandocbounded{\includegraphics[keepaspectratio]{04-resultados-curso_files/figure-pdf/fig-n-carrera-psico-1.pdf}}

}

\caption{\label{fig-n-carrera-psico}Elaboración propia (2026)}

\end{figure}%

\subsubsection{Antropología}

\begin{figure}

\centering{

\pandocbounded{\includegraphics[keepaspectratio]{04-resultados-curso_files/figure-pdf/fig-n-ant-1.pdf}}

}

\caption{\label{fig-n-ant}Elaboración propia (2026)}

\end{figure}%

\subsubsection{Trabajo Social}

\begin{figure}

\centering{

\pandocbounded{\includegraphics[keepaspectratio]{04-resultados-curso_files/figure-pdf/fig-n-ts-1.pdf}}

}

\caption{\label{fig-n-ts}Elaboración propia (2026)}

\end{figure}%

\subsubsection{Educación Parvularia}

\begin{figure}

\centering{

\pandocbounded{\includegraphics[keepaspectratio]{04-resultados-curso_files/figure-pdf/fig-n-ed-1.pdf}}

}

\caption{\label{fig-n-ed}Elaboración propia (2026)}

\end{figure}%

Este grupo de gráficos desagrega la relación entre tamaño del curso y
las notas por carrera. Se observa un patrón general donde las
diferencias son pequeñas y poco sistemáticas en todas las unidades. En
Sociología se aprecia una leve tendencia decreciente: los cursos más
pequeños (≤20) tienen medianas ligeramente superiores
(\textasciitilde6.1) y los más grandes (≥81) muestran menores
calificaciones y, en algunos casos, mayor dispersión. En Psicología el
patrón es similar pero atenuado, con medianas más altas para los cursos
de tamaño reducido y menores calificaciones acompañadas de mayor
dispersión para cursos más numerosos . En Antropología no existe ninguna
tendencia clara: todos los tamaños muestran medianas similares
(\textasciitilde6.0) y dispersión homogénea con outliers distribuidos en
varios grupos. Por el lado de Trabajo Social los cursos más grandes
(61--80) presentan medianas más bajas (\textasciitilde5.8) respecto a
los más pequeños, pero cuenta sólo con cuatro categorías. Educación
Parvularia muestra un comportamiento uniforme, con medianas estables en
torno a 6.0--6.1 independientemente del tamaño. En conjunto, el tamaño
del curso presenta ciertas diferencias al analizar por carrera, pero en
general no se encuentran diferencias sustantivas.

\subsubsection{Carreras en el tiempo}\label{carreras-en-el-tiempo-2}

\subsubsection{Sociología}

\begin{figure}

\centering{

\pandocbounded{\includegraphics[keepaspectratio]{04-resultados-curso_files/figure-pdf/fig-n-socio-t-1.pdf}}

}

\caption{\label{fig-n-socio-t}Elaboración propia (2026)}

\end{figure}%

\subsubsection{Psicología}

\begin{figure}

\centering{

\pandocbounded{\includegraphics[keepaspectratio]{04-resultados-curso_files/figure-pdf/fig-n-psico-t-1.pdf}}

}

\caption{\label{fig-n-psico-t}Elaboración propia (2026)}

\end{figure}%

\subsubsection{Antropología}

\begin{figure}

\centering{

\pandocbounded{\includegraphics[keepaspectratio]{04-resultados-curso_files/figure-pdf/fig-n-ant-t-1.pdf}}

}

\caption{\label{fig-n-ant-t}Elaboración propia (2026)}

\end{figure}%

\subsubsection{Trabajo Social}

\begin{figure}

\centering{

\pandocbounded{\includegraphics[keepaspectratio]{04-resultados-curso_files/figure-pdf/fig-n-ts-t-1.pdf}}

}

\caption{\label{fig-n-ts-t}Elaboración propia (2026)}

\end{figure}%

\subsubsection{Educación Parvularia}

\begin{figure}

\centering{

\pandocbounded{\includegraphics[keepaspectratio]{04-resultados-curso_files/figure-pdf/fig-n-ed-t-1.pdf}}

}

\caption{\label{fig-n-ed-t}Elaboración propia (2026)}

\end{figure}%

Este grupo de gráficos presenta la trayectoria de notas según tamaño del
curso por carrera, y el hallazgo central es que no existe un patrón
sistemático que pueda generalizarse entre carreras. En Sociología, los
cursos más pequeños (≤20) muestran una ventaja leve y creciente hacia
2025, mientras los más grandes (≥81) exhiben la trayectoria más baja e
inestable, aunque todas las líneas se entrelazan frecuentemente. En
Psicología, las líneas son altamente volátiles y se cruzan de forma
irregular, sin que ningún tamaño mantenga una ventaja sostenida,
destacando solo una caída pronunciada del grupo 61--80 en 2024. En
Antropología, todas las categorías oscilan en un rango estrecho
(\textasciitilde5.8--6.1) con trayectorias entrelazadas y sin
diferencias apreciables. Por su parte, en Trabajo Social, el grupo
21--40 muestra una caída abrupta hacia 2024--2025 que lo distingue
negativamente del resto, mientras que los cursos más grandes (61-80)
experimentan un alza notoria el 2025. Educación Parvularia presenta el
comportamiento más estable, con la mayoría de los grupos convergiendo en
torno a 6.0--6.4 desde 2020 en adelante. Es llamativo que Trabajo Social
y Educación Parvularia no tengan la categoría más alta de tamaños de
curso, pero esto se debe a que las carreras no admiten tales cantidades
de estudiantes por año, resultando en un límite de la planificación de
matrículas En síntesis, desagregando por carrera, el tamaño del curso no
explica de forma sistemática las diferencias en rendimiento, confirmando
la conclusión del análisis transversal previo.

\bookmarksetup{startatroot}

\chapter{Panel interactivo Cursos
FACSO}\label{panel-interactivo-cursos-facso}

\bookmarksetup{startatroot}

\chapter{Conclusiones}\label{conclusiones}

El Estudio de Notas de la Facultad de Ciencias Sociales ha entregado
luces acerca de cómo se distribuyen las calificaciones, enfocándose en
los estudiantes y cursos a partir de un nivel agregado de facultad, como
también en las dinámicas de cada carrera. Así, los resultados observados
permiten responder las preguntas que fueron planteadas en la
introducción del reporte.

\textbf{¿Existe la inflación de notas en la Facultad de Ciencias
Sociales?}

A partir de los elementos que circunscriben el análisis, se encontró
evidencia que puede sustentar la hipótesis de inflación de notas. El
alza pronunciada de las notas en los años 2020 y 2021 en comparación de
años prepandémicos y su persistencia hasta el año 2025 demuestran que
FACSO atraviesa un proceso que actualmente está presente en la educación
superior a nivel global.

\textbf{¿Las notas permiten discriminar grupos de estudiantes en función
del rendimiento?}

Las calificaciones dificultan realizar una distinción entre estudiantes
que sobresalen académicamente de aquellos que necesitan apoyo para
cumplir con los aprendizajes esperados debido a la baja variabilidad de
las notas que obtienen los estudiantes. Incluso, las calificaciones han
tendido a homogeneizarse con el paso del tiempo, complicando aún más la
tarea de diferenciar grupos de estudiantes a partir de su rendimiento
académico. Esto es sumamente preocupante, ya que no se cuenta con la
información necesaria para gestionar e implementar políticas enfocadas
en estudiantes que requieren un apoyo adicional en su trayectoria
académica, con el objetivo de reducir las brechas educativas de base.

\textbf{¿Qué tan alto es el nivel de exigencia en cuánto a la aprobación
de asignaturas?}

Los niveles de exigencia de los cursos en FACSO pueden entenderse como
bajos, puesto que entre 2018 y 2025 el promedio de aprobación de los
cursos fue de un 97\%. Esto significa que reprobar un curso es una
experiencia anómala, donde no hay mayores dificultades en acreditar la
obtención de los conocimientos esperados de un ramo.

\textbf{¿Qué factores asociados dan cuenta de las notas, el grado de
discriminación y la tasa de aprobación?}

En primer lugar, factores socioeconómicos como el tipo de colegio y el
nivel socioeconómico demostraron ser determinantes en las calificaciones
de los estudiantes. En segundo lugar, el género evidenció tener un rol
explicativo importante, donde las mujeres tienen un rendimiento más alto
que su contraparte. En tercer lugar, variables institucionales como el
tipo de curso y el tamaño del curso se confirmaron como factores
asociados a las calificaciones.

Considerando el conjunto de estos hallazgos se asoma como un imperativo
desarrollar una discusión seria sobre el panorama actual de la dimensión
formativa-evaluativa de la facultad: en un contexto de inflación de
notas, el hecho que las evaluaciones no hagan posible diferenciar a
grupos de estudiantes puede tener consecuencias graves en la población
estudiantil, como sentimientos de injusticia en la calificación
académica. Además, se podría hipotetizar una carencia formativa de los
estudiantes, pero precisamente la alta homogeneidad en las
calificaciones no permite identificar si tales carencias existen.

En el marco de la incorporación de normativas que promulgan la no
discriminación y la equidad, hay que tener suma precaución con que estos
principios no se confundan con la capacidad de capturar las diferencias
sociales y académicas que, como institución de excelencia, la
Universidad de Chile debe ser capaz de observar y abordar
apropiadamente. Por ello, es un deber institucional generar estrategias
que puedan comprender cómo se plantean y despliegan las evaluaciones a
estudiantes. Con todo, este estudio representa un diagnóstico de las
dinámicas en las calificaciones académicas, y puede significar un insumo
sumamente valioso para las autoridades al momento de promover de manera
enfocada los recursos necesarios para pulir los sistemas de evaluación
y, en última instancia, acortar las brechas en el logro académico de los
estudiantes.

En el futuro, este estudio podría entregar más pistas sobre las
calificaciones en la facultad a partir de la incorporación de
información nueva. En este contexto, contar con las evaluaciones
docentes y los cursos impartidos a lo largo del tiempo podrían aportar a
caracterizar con mayor profundidad qué cursos cuentan con mayores tasas
de aprobación, y cómo eso se relaciona con el juicio que realizan los
estudiantes de los cursos. Asimismo, sería valioso realizar un estudio
comparativo de los criterios de evaluación que declaran los cursos de
cada carrera en sus programas, puesto que se podrían mapear los modos de
capturar los aprendizajes, la asignación de métricas y los umbrales que
determinan el éxito del proceso. Como última propuesta, replicar el este
estudio en otras facultades significaría un avance en la comprensión del
escenario formativo-evaluativo a nivel Universidad de Chile, ya que
permitiría comparar las diferencias entre facultades que forman parte de
un mismo ecosistema académico.

\bookmarksetup{startatroot}

\chapter{Referencias}\label{referencias}

\phantomsection\label{refs}
\begin{CSLReferences}{1}{0}
\bibitem[\citeproctext]{ref-achen_what_2009}
Achen, A. C., \& Courant, P. N. (2009). What Are Grades Made Of?
\emph{Journal of Economic Perspectives}, \emph{23}(3), 77-92.
\url{https://doi.org/10.1257/jep.23.3.77}

\bibitem[\citeproctext]{ref-hernandez-julian_measuring_2016}
Hernández-Julián, R., \& Looney, A. (2016). Measuring Inflation in
Grades: An Application of Price Indexing to Undergraduate Grades.
\emph{Economics of Education Review}, \emph{55}, 220-232.
\url{https://doi.org/10.1016/j.econedurev.2016.11.001}

\bibitem[\citeproctext]{ref-jephcote_grade_2021}
Jephcote, C., Medland, E., \& Lygo-Baker, S. (2021). Grade Inflation
versus Grade Improvement: Are Our Students Getting More Intelligent?
\emph{Assessment \& Evaluation in Higher Education}, \emph{46}(4),
547-571. \url{https://doi.org/10.1080/02602938.2020.1795617}

\bibitem[\citeproctext]{ref-jewell_whose_2011}
Jewell, R. T., McPherson, M. A., \& Tieslau, M. A. (2011). Whose Fault
Is It? Assigning Blame for Grade Inflation in Higher Education.
\emph{Applied Economics}, \emph{45}(9), 1185-1200.
\url{https://doi.org/10.1080/00036846.2011.621884}

\bibitem[\citeproctext]{ref-lipnevich_what_2020}
Lipnevich, A. A., Guskey, T. R., Murano, D. M., \& Smith, J. K. (2020).
What Do Grades Mean? Variation in Grading Criteria in American College
and University Courses. \emph{Assessment in Education: Principles,
Policy \& Practice}, \emph{27}(5), 480-500.
\url{https://doi.org/10.1080/0969594X.2020.1799190}

\bibitem[\citeproctext]{ref-lynch_learning_2017}
Lynch, R., \& Hennessy, J. (2017). Learning to Earn? The Role of
Performance Grades in Higher Education. \emph{Studies in Higher
Education}, \emph{42}(9), 1750-1763.
\url{https://doi.org/10.1080/03075079.2015.1124850}

\bibitem[\citeproctext]{ref-schorr_addressing_2025}
Schorr, C. (2025). Addressing the Grade Inflation Collective Action
Problem \textbar{} U.S. Department of Education. En \emph{U. S.
Department of Education}.
http://www.ed.gov/about/homeroom-blog/addressing-grade-inflation-collective-action-problem.

\bibitem[\citeproctext]{ref-universidaddechile_modelo_2024}
Universidad de Chile. (2024). \emph{Modelo educativo de la Universidad
de Chile} (Universidad de Chile).

\end{CSLReferences}

\cleardoublepage
\phantomsection
\addcontentsline{toc}{part}{Apéndices}
\appendix

\chapter{Anexos estudiantes}\label{anexos-estudiantes}

\section{Cantidad de notas según categoría de
agrupación}\label{cantidad-de-notas-seguxfan-categoruxeda-de-agrupaciuxf3n}

\begin{longtable}[t]{lrrr}

\caption{\label{tbl-group-notas}Cantidad de notas según categoría
(2020-2024)}

\tabularnewline

\toprule
\cellcolor[HTML]{952626}{\textcolor{white}{\textbf{Categoría nota}}} & \cellcolor[HTML]{952626}{\textcolor{white}{\textbf{Frec. absoluta}}} & \cellcolor[HTML]{952626}{\textcolor{white}{\textbf{Frec. relativa}}} & \cellcolor[HTML]{952626}{\textcolor{white}{\textbf{Frec. acumulada}}}\\
\midrule
<= 7,0 y >= 6,5 & 146 & 0.07 & 0.07\\
<= 6,4 y >= 6,0 & 926 & 0.47 & 0.54\\
<= 5,9 y >= 5,5 & 628 & 0.32 & 0.86\\
<= 5,4 y >= 5,0 & 180 & 0.09 & 0.95\\
<= 4,9 & 101 & 0.05 & 1.00\\
\addlinespace
\cellcolor[HTML]{f2f2f2}{\textbf{Total}} & \cellcolor[HTML]{f2f2f2}{\textbf{1981}} & \cellcolor[HTML]{f2f2f2}{\textbf{1.00}} & \cellcolor[HTML]{f2f2f2}{\textbf{1.00}}\\
\bottomrule

\end{longtable}

\section{Notas promedio de cada carrera por
cohorte}\label{notas-promedio-de-cada-carrera-por-cohorte}

\begin{table}

\caption{\label{tbl-desc-carreras}Estadísticas descriptivas del promedio
por carrera y cohorte}

\centering{

\centering
\resizebox{\ifdim\width>\linewidth\linewidth\else\width\fi}{!}{
\begin{tabular}{>{\raggedright\arraybackslash}p{4cm}|c|c|c|c|l}
\hline
\textbf{carrera} & \textbf{2020} & \textbf{2021} & \textbf{2022} & \textbf{2023} & \textbf{2024}\\
\hline
\textbf{Antropología} & 6 (DE = 0.4) [4.8–6.7] & 5.9 (DE = 0.4) [4–6.7] & 5.8 (DE = 0.6) [2.6–6.7] & 5.8 (DE = 0.4) [4.4–6.4] & 5.8 (DE = 0.4) [3.5–6.5]\\
\cline{1-6}
\textbf{Educación Parvularia} & 6.4 (DE = 0.2) [5.8–6.7] & 6.4 (DE = 0.2) [5.9–6.8] & 6.2 (DE = 0.3) [5.2–6.8] & 6.2 (DE = 0.4) [4.2–6.7] & 5.9 (DE = 1.2) [1.5–6.8]\\
\cline{1-6}
\textbf{Psicología} & 6.4 (DE = 0.3) [5–6.8] & 6.2 (DE = 0.3) [5.1–6.8] & 6.1 (DE = 0.5) [2–6.5] & 5.8 (DE = 0.8) [2.6–6.6] & 5.6 (DE = 0.6) [2.8–6.5]\\
\cline{1-6}
\textbf{Sociología} & 6.1 (DE = 0.5) [2.6–6.6] & 5.9 (DE = 0.7) [2.6–6.6] & 5.9 (DE = 0.4) [4.5–6.5] & 5.8 (DE = 0.5) [3.7–6.7] & 5.7 (DE = 0.4) [3.4–6.5]\\
\cline{1-6}
\textbf{Trabajo Social} & 6 (DE = 0.3) [5.3–6.5] & 6 (DE = 0.4) [4.2–6.5] & 5.9 (DE = 0.4) [4.5–6.4] & 5.6 (DE = 0.7) [2.6–6.3] & 5.4 (DE = 1.1) [1.6–6.4]\\
\hline
\end{tabular}}

}

\end{table}%

\section{Género por cohortes}\label{guxe9nero-por-cohortes}

\begin{figure}

\caption{\label{fig-plotly-gen}Elaboración propia (2026)}

\centering{

\pandocbounded{\includegraphics[keepaspectratio]{apendices/A-anexos-estudiantes_files/figure-pdf/fig-plotly-gen-1.pdf}}

}

\end{figure}%

\section{Tipo colegio por cohortes}\label{tipo-colegio-por-cohortes}

\begin{figure}

\caption{\label{fig-time-col}Elaboración propia (2026)}

\centering{

\pandocbounded{\includegraphics[keepaspectratio]{apendices/A-anexos-estudiantes_files/figure-pdf/fig-time-col-1.pdf}}

}

\end{figure}%

\#\# Nivel socioeconómico por cohortes

\begin{figure}

\caption{\label{fig-plotly-nse}Elaboración propia (2026)}

\centering{

\pandocbounded{\includegraphics[keepaspectratio]{apendices/A-anexos-estudiantes_files/figure-pdf/fig-plotly-nse-1.pdf}}

}

\end{figure}%

\section{Percentil prueba admisión universitaria por
cohortes}\label{percentil-prueba-admisiuxf3n-universitaria-por-cohortes}

\begin{figure}

\caption{\label{fig-scat-pct}Elaboración propia (2026)}

\centering{

\pandocbounded{\includegraphics[keepaspectratio]{apendices/A-anexos-estudiantes_files/figure-pdf/fig-scat-pct-1.pdf}}

}

\end{figure}%

\section{Percentil prueba admisión universitaria por
carrera}\label{percentil-prueba-admisiuxf3n-universitaria-por-carrera}

\begin{figure}

\centering{

\pandocbounded{\includegraphics[keepaspectratio]{apendices/A-anexos-estudiantes_files/figure-pdf/fig-admision-1.pdf}}

}

\caption{\label{fig-admision}Elaboración propia en base a datos de
UCampus (2026)}

\end{figure}%

\begin{figure}

\caption{\label{fig-scat-socio}Elaboración propia en base a datos de
UCampus (2026)}

\centering{

\pandocbounded{\includegraphics[keepaspectratio]{apendices/A-anexos-estudiantes_files/figure-pdf/fig-scat-socio-1.pdf}}

}

\end{figure}%

\begin{figure}

\caption{\label{fig-scat-psico}Elaboración propia en base a datos de
UCampus (2026)}

\centering{

\pandocbounded{\includegraphics[keepaspectratio]{apendices/A-anexos-estudiantes_files/figure-pdf/fig-scat-psico-1.pdf}}

}

\end{figure}%

\begin{figure}

\caption{\label{fig-scat-antropo}Elaboración propia en base a datos de
UCampus (2026)}

\centering{

\pandocbounded{\includegraphics[keepaspectratio]{apendices/A-anexos-estudiantes_files/figure-pdf/fig-scat-antropo-1.pdf}}

}

\end{figure}%

\begin{figure}

\caption{\label{fig-scat-educa}Elaboración propia en base a datos de
UCampus (2026)}

\centering{

\pandocbounded{\includegraphics[keepaspectratio]{apendices/A-anexos-estudiantes_files/figure-pdf/fig-scat-educa-1.pdf}}

}

\end{figure}%

\begin{figure}

\caption{\label{fig-scat-trabajo}Elaboración propia en base a datos de
UCampus (2026)}

\centering{

\pandocbounded{\includegraphics[keepaspectratio]{apendices/A-anexos-estudiantes_files/figure-pdf/fig-scat-trabajo-1.pdf}}

}

\end{figure}%

\section{Comparación del top 40}\label{comparaciuxf3n-del-top-40}

\begin{figure}

\caption{\label{fig-socio-top}Comparación de las mejores 40 notas de
Sociología por cohorte}

\centering{

\pandocbounded{\includegraphics[keepaspectratio]{apendices/A-anexos-estudiantes_files/figure-pdf/fig-socio-top-1.pdf}}

}

\end{figure}%

\begin{figure}

\caption{\label{fig-psi-top}Comparación de las mejores 40 notas de
Psicología por cohorte}

\centering{

\pandocbounded{\includegraphics[keepaspectratio]{apendices/A-anexos-estudiantes_files/figure-pdf/fig-psi-top-1.pdf}}

}

\end{figure}%

\begin{figure}

\caption{\label{fig-trabajo-top}Comparación de las mejores 40 notas de
Trabajo Social por cohorte}

\centering{

\pandocbounded{\includegraphics[keepaspectratio]{apendices/A-anexos-estudiantes_files/figure-pdf/fig-trabajo-top-1.pdf}}

}

\end{figure}%

\begin{figure}

\caption{\label{fig-antropo-top}Comparación de las mejores 40 notas de
Antropología por cohorte}

\centering{

\pandocbounded{\includegraphics[keepaspectratio]{apendices/A-anexos-estudiantes_files/figure-pdf/fig-antropo-top-1.pdf}}

}

\end{figure}%

\begin{figure}

\caption{\label{fig-educa-top}Comparación de las mejores 40 notas de
Educación Parvularia por cohorte}

\centering{

\pandocbounded{\includegraphics[keepaspectratio]{apendices/A-anexos-estudiantes_files/figure-pdf/fig-educa-top-1.pdf}}

}

\end{figure}%

\section{Comparación del bottom 40}\label{comparaciuxf3n-del-bottom-40}

\begin{figure}

\caption{\label{fig-socio-bottom}Comparación de las 40 notas más bajas
de Sociología por cohorte}

\centering{

\pandocbounded{\includegraphics[keepaspectratio]{apendices/A-anexos-estudiantes_files/figure-pdf/fig-socio-bottom-1.pdf}}

}

\end{figure}%

\begin{figure}

\caption{\label{fig-psi-bottom}Comparación de las 40 notas más bajas de
Psicología por cohorte}

\centering{

\pandocbounded{\includegraphics[keepaspectratio]{apendices/A-anexos-estudiantes_files/figure-pdf/fig-psi-bottom-1.pdf}}

}

\end{figure}%

\begin{figure}

\caption{\label{fig-trabajo-bottom}Comparación de las 40 notas más bajas
de Trabajo Social por cohorte}

\centering{

\pandocbounded{\includegraphics[keepaspectratio]{apendices/A-anexos-estudiantes_files/figure-pdf/fig-trabajo-bottom-1.pdf}}

}

\end{figure}%

\begin{figure}

\caption{\label{fig-antropo-bottom}Comparación de las 40 notas más bajas
de Antropología por cohorte}

\centering{

\pandocbounded{\includegraphics[keepaspectratio]{apendices/A-anexos-estudiantes_files/figure-pdf/fig-antropo-bottom-1.pdf}}

}

\end{figure}%

\begin{figure}

\caption{\label{fig-educa-bottom}Comparación de las 40 notas más bajas
de Educación Parvularia por cohorte}

\centering{

\pandocbounded{\includegraphics[keepaspectratio]{apendices/A-anexos-estudiantes_files/figure-pdf/fig-educa-bottom-1.pdf}}

}

\end{figure}%

\section{Análisis de notas a través del
tiempo}\label{anuxe1lisis-de-notas-a-travuxe9s-del-tiempo}

Se adjunta el código con el cual se construyeron las bases de los cursos
de primer año. Este código fue ejecutado en \texttt{proc-bbdd.qmd}, el
cual se encuentra en la carpeta \texttt{proc}.

\section{Promedio de la Facultad a través del
tiempo}\label{promedio-de-la-facultad-a-travuxe9s-del-tiempo}

\begin{figure}

\caption{\label{fig-año-facso}Promedio de notas por año 2021 - 2024
FACSO}

\centering{

\pandocbounded{\includegraphics[keepaspectratio]{apendices/A-anexos-estudiantes_files/figure-pdf/fig-año-facso-1.pdf}}

}

\end{figure}%

\section{Promedios de carrera a través del
tiempo}\label{promedios-de-carrera-a-travuxe9s-del-tiempo}

\begin{figure}

\caption{\label{fig-año-socio}Promedio de notas por año 2021 - 2024
Sociología}

\centering{

\pandocbounded{\includegraphics[keepaspectratio]{apendices/A-anexos-estudiantes_files/figure-pdf/fig-año-socio-1.pdf}}

}

\end{figure}%

\begin{figure}

\caption{\label{fig-año-psico}Promedio de notas por año 2021 - 2024
Psicología}

\centering{

\pandocbounded{\includegraphics[keepaspectratio]{apendices/A-anexos-estudiantes_files/figure-pdf/fig-año-psico-1.pdf}}

}

\end{figure}%

\begin{figure}

\caption{\label{fig-año-antropo}Promedio de notas por año 2021 - 2024
Antropología}

\centering{

\pandocbounded{\includegraphics[keepaspectratio]{apendices/A-anexos-estudiantes_files/figure-pdf/fig-año-antropo-1.pdf}}

}

\end{figure}%

\begin{figure}

\caption{\label{fig-año-educa}Promedio de notas por año 2021 - 2024
Educación Parvularia}

\centering{

\pandocbounded{\includegraphics[keepaspectratio]{apendices/A-anexos-estudiantes_files/figure-pdf/fig-año-educa-1.pdf}}

}

\end{figure}%

\begin{figure}

\caption{\label{fig-año-trabajo}Promedio de notas por año 2021 - 2024
Trabajo Social}

\centering{

\pandocbounded{\includegraphics[keepaspectratio]{apendices/A-anexos-estudiantes_files/figure-pdf/fig-año-trabajo-1.pdf}}

}

\end{figure}%

\pandocbounded{\includegraphics[keepaspectratio]{apendices/A-anexos-estudiantes_files/figure-pdf/unnamed-chunk-16-1.pdf}}

\begin{figure}

\caption{\label{fig-facso-full}Asociación entre Promedio por defecto y
Promedio calculado - FACSO}

\centering{

\pandocbounded{\includegraphics[keepaspectratio]{apendices/A-anexos-estudiantes_files/figure-pdf/fig-facso-full-1.pdf}}

}

\end{figure}%

\begin{figure}

\caption{\label{fig-socio-full}Asociación entre Promedio bruto y
Promedio calculado - Sociología}

\centering{

\pandocbounded{\includegraphics[keepaspectratio]{apendices/A-anexos-estudiantes_files/figure-pdf/fig-socio-full-1.pdf}}

}

\end{figure}%

\begin{figure}

\caption{\label{fig-psico-full}Asociación entre Promedio bruto y
Promedio calculado - Psicología}

\centering{

\pandocbounded{\includegraphics[keepaspectratio]{apendices/A-anexos-estudiantes_files/figure-pdf/fig-psico-full-1.pdf}}

}

\end{figure}%

\begin{figure}

\caption{\label{fig-trabajo-full}Asociación entre Promedio bruto y
Promedio calculado - Trabajo social}

\centering{

\pandocbounded{\includegraphics[keepaspectratio]{apendices/A-anexos-estudiantes_files/figure-pdf/fig-trabajo-full-1.pdf}}

}

\end{figure}%

\begin{figure}

\caption{\label{fig-educa-full}Asociación entre Promedio bruto y
Promedio calculado - Educación Parvularia}

\centering{

\pandocbounded{\includegraphics[keepaspectratio]{apendices/A-anexos-estudiantes_files/figure-pdf/fig-educa-full-1.pdf}}

}

\end{figure}%

\begin{figure}

\caption{\label{fig-antropo-full}Asociación entre Promedio bruto y
Promedio calculado - Antropología Arqueología}

\centering{

\pandocbounded{\includegraphics[keepaspectratio]{apendices/A-anexos-estudiantes_files/figure-pdf/fig-antropo-full-1.pdf}}

}

\end{figure}%

\chapter{Anexos cursos}\label{anexos-cursos}

\begin{longtable}[t]{lrrrrr>{}r}

\caption{\label{tbl-n-cursos}Cantidad de cursos por año y carrera
(2018-2025)}

\tabularnewline

\toprule
\cellcolor[HTML]{952626}{\textcolor{white}{\textbf{Año}}} & \cellcolor[HTML]{952626}{\textcolor{white}{\textbf{Antropología}}} & \cellcolor[HTML]{952626}{\textcolor{white}{\textbf{Educación Parvularia}}} & \cellcolor[HTML]{952626}{\textcolor{white}{\textbf{Psicología}}} & \cellcolor[HTML]{952626}{\textcolor{white}{\textbf{Sociología}}} & \cellcolor[HTML]{952626}{\textcolor{white}{\textbf{Trabajo Social}}} & \cellcolor[HTML]{952626}{\textcolor{white}{\textbf{Total}}}\\
\midrule
2018 & 12 & 12 & 36 & 33 & 25 & \textbf{118}\\
2019 & 29 & 25 & 51 & 60 & 38 & \textbf{203}\\
2020 & 93 & 37 & 88 & 92 & 54 & \textbf{364}\\
2021 & 153 & 68 & 102 & 126 & 67 & \textbf{516}\\
2022 & 219 & 103 & 109 & 195 & 69 & \textbf{695}\\
\addlinespace
2023 & 225 & 103 & 135 & 187 & 70 & \textbf{720}\\
2024 & 218 & 132 & 423 & 178 & 94 & \textbf{1045}\\
2025 & 220 & 147 & 298 & 193 & 140 & \textbf{998}\\
\cellcolor[HTML]{f2f2f2}{\textbf{Total}} & \cellcolor[HTML]{f2f2f2}{\textbf{1169}} & \cellcolor[HTML]{f2f2f2}{\textbf{627}} & \cellcolor[HTML]{f2f2f2}{\textbf{1242}} & \cellcolor[HTML]{f2f2f2}{\textbf{1064}} & \cellcolor[HTML]{f2f2f2}{\textbf{557}} & \textbf{\cellcolor[HTML]{f2f2f2}{\textbf{4659}}}\\
\bottomrule

\end{longtable}

\begin{longtable}[t]{lrrrr}

\caption{\label{tbl-n-notas}Descriptivos según tamaño del curso
(2018-2025)}

\tabularnewline

\toprule
\cellcolor[HTML]{952626}{\textcolor{white}{\textbf{Tamaño del curso}}} & \cellcolor[HTML]{952626}{\textcolor{white}{\textbf{Frec. absoluta}}} & \cellcolor[HTML]{952626}{\textcolor{white}{\textbf{Frec. relativa}}} & \cellcolor[HTML]{952626}{\textcolor{white}{\textbf{Nota promedio}}} & \cellcolor[HTML]{952626}{\textcolor{white}{\textbf{Nota prom. acumulada}}}\\
\midrule
1 estudiante & 689 & 0.15 & 6.36 & 6.36\\
> 1 y <= 5 & 1346 & 0.29 & 6.34 & 6.34\\
> 5 y <= 10 & 387 & 0.08 & 6.22 & 6.32\\
> 10 y <= 20 & 435 & 0.09 & 6.08 & 6.29\\
> 20 y <= 40 & 783 & 0.17 & 6.11 & 6.25\\
\addlinespace
> 40 & 1024 & 0.22 & 5.90 & 6.17\\
\cellcolor[HTML]{f2f2f2}{\textbf{Total}} & \cellcolor[HTML]{f2f2f2}{\textbf{4664}} & \cellcolor[HTML]{f2f2f2}{\textbf{1.00}} & \cellcolor[HTML]{f2f2f2}{\textbf{37.00}} & \cellcolor[HTML]{f2f2f2}{\textbf{6.36}}\\
\bottomrule

\end{longtable}

\begin{longtable}[t]{lrrr}

\caption{\label{tbl-cat-notas}Cantidad de cursos según nota promedio
(2018-2025)}

\tabularnewline

\toprule
\cellcolor[HTML]{952626}{\textcolor{white}{\textbf{Categoría nota}}} & \cellcolor[HTML]{952626}{\textcolor{white}{\textbf{Frec. absoluta}}} & \cellcolor[HTML]{952626}{\textcolor{white}{\textbf{Frec. relativa}}} & \cellcolor[HTML]{952626}{\textcolor{white}{\textbf{Frec. acumulada}}}\\
\midrule
< 4,0 & 68 & 0.01 & 0.01\\
>= 4,0 y < 5,0 & 168 & 0.04 & 0.05\\
>= 5,0 y < 6,0 & 1268 & 0.27 & 0.32\\
>= 6,0 y < 6,5 & 1402 & 0.30 & 0.62\\
\cellcolor[HTML]{f2f2f2}{\textbf{>= 6,5}} & \cellcolor[HTML]{f2f2f2}{\textbf{1758}} & \cellcolor[HTML]{f2f2f2}{\textbf{0.38}} & \cellcolor[HTML]{f2f2f2}{\textbf{1.00}}\\
\bottomrule

\end{longtable}

\begin{longtable}[t]{lrrr}

\caption{\label{tbl-cat-aprob}Cantidad de cursos según tasa de
aprobación (2018-2025)}

\tabularnewline

\toprule
\cellcolor[HTML]{952626}{\textcolor{white}{\textbf{\% aprobación}}} & \cellcolor[HTML]{952626}{\textcolor{white}{\textbf{Frec. absoluta}}} & \cellcolor[HTML]{952626}{\textcolor{white}{\textbf{Frec. relativa}}} & \cellcolor[HTML]{952626}{\textcolor{white}{\textbf{Frec. acumulada}}}\\
\midrule
< 90 & 329 & 0.07 & 0.07\\
>= 90 y < 92 & 93 & 0.02 & 0.09\\
>= 92 y < 94 & 136 & 0.03 & 0.12\\
>= 94 y < 96 & 175 & 0.04 & 0.16\\
>= 96 y < 98 & 310 & 0.07 & 0.22\\
\addlinespace
>= 98 y < 100 & 172 & 0.04 & 0.26\\
\cellcolor[HTML]{f2f2f2}{\textbf{100}} & \cellcolor[HTML]{f2f2f2}{\textbf{3449}} & \cellcolor[HTML]{f2f2f2}{\textbf{0.74}} & \cellcolor[HTML]{f2f2f2}{\textbf{1.00}}\\
\bottomrule

\end{longtable}

\begin{longtable}[t]{lllrr}

\caption{\label{tbl-conteo-notas}Listado de cursos con notas promedio
menores a 4.0}

\tabularnewline

\toprule
\cellcolor[HTML]{952626}{\textcolor{white}{\textbf{Año}}} & \cellcolor[HTML]{952626}{\textcolor{white}{\textbf{Carrera}}} & \cellcolor[HTML]{952626}{\textcolor{white}{\textbf{Curso}}} & \cellcolor[HTML]{952626}{\textcolor{white}{\textbf{Nota promedio}}} & \cellcolor[HTML]{952626}{\textcolor{white}{\textbf{N estudiantes}}}\\
\midrule
2018 & Psicología & Estadística 2 & 1 & 3.30\\
2018 & Sociología & Economía & 2 & 2.75\\
2018 & Sociología & Teorías Sociológicas Contemporáneas & 5 & 3.02\\
2018 & Trabajo Social & Planificación y Gestión Social & 1 & 1.00\\
2019 & Sociología & Electivo de Especialidad - Sociología Urbana & 1 & 1.40\\
\addlinespace
2019 & Sociología & Desigualdades y Estratificación Social & 3 & 3.17\\
2020 & Trabajo Social & Núcleo Temático Electivo I + D III: Estudios Interdisciplinarios en Trabajo Social & 3 & 3.67\\
2021 & Educación Parvularia & Seminario Temático de Investigación & 2 & 3.75\\
2021 & Educación Parvularia & Seminario Temático de Investigación & 2 & 3.85\\
2022 & Antropología & Taller de Memoria I: Mención Social & 1 & 1.00\\
\addlinespace
2022 & Antropología & Taller de Investigación II - Arqueología & 1 & 1.00\\
2022 & Antropología & Taller de Investigación I - Social & 1 & 1.00\\
2022 & Antropología & Taller de Investigación I - Social & 5 & 3.20\\
2022 & Antropología & Histología & 16 & 3.76\\
2022 & Antropología & Taller de Investigación I - Social & 1 & 3.90\\
\addlinespace
2022 & Educación Parvularia & Seminario Título II & 1 & 1.00\\
2022 & Educación Parvularia & Seminario Temático de Investigación & 2 & 3.05\\
2022 & Educación Parvularia & Seminario Título II & 1 & 3.60\\
2023 & Antropología & Taller de Memoria I: Mención Social & 1 & 1.00\\
2023 & Antropología & Taller de Investigación I - Social & 1 & \vphantom{1} 1.00\\
\addlinespace
2023 & Antropología & Taller de Investigación I - Social & 1 & 1.00\\
2023 & Antropología & Taller de Memoria II: Mención Social & 1 & 1.00\\
2023 & Antropología & Taller de Memoria I: Mención Social & 1 & 3.00\\
2023 & Antropología & Taller de Integración: Ciencia y Tecnología & 8 & 3.09\\
2023 & Antropología & Taller de Investigación II - Social & 2 & 3.50\\
\addlinespace
2023 & Antropología & Taller de Investigación I - Social & 2 & 3.85\\
2023 & Antropología & Taller de Memoria I Mención Arqueología & 2 & 3.90\\
2023 & Antropología & Taller de Investigación I - Social & 4 & 3.90\\
2023 & Educación Parvularia & Seminario Título II & 1 & 3.60\\
2023 & Sociología & Seminario de Título II & 2 & 3.00\\
\addlinespace
2023 & Sociología & Seminario de Título II & 1 & \vphantom{1} 3.50\\
2023 & Sociología & Seminario de Título II & 1 & 3.50\\
2023 & Sociología & Seminario de Título II & 2 & 3.75\\
2023 & Sociología & Seminario de Grado & 3 & 3.77\\
2024 & Antropología & Taller de Memoria I: Mención Social & 1 & 1.00\\
\addlinespace
2024 & Antropología & Taller de Investigación II - Arqueología & 1 & 1.00\\
2024 & Antropología & Taller de Investigación II - Social & 1 & \vphantom{1} 3.00\\
2024 & Antropología & Taller de Investigación II - Social & 1 & 3.00\\
2024 & Antropología & Taller de Investigación I - Social & 2 & 3.40\\
2024 & Antropología & Taller de Investigación II - Física & 2 & 3.55\\
\addlinespace
2024 & Antropología & Taller de Investigación II - Arqueología & 2 & 3.65\\
2024 & Antropología & Taller de Investigación I - Social & 3 & 3.77\\
2024 & Antropología & Taller de Investigación II - Arqueología & 4 & 3.80\\
2024 & Antropología & Taller de Memoria I Mención Arqueología & 5 & 3.80\\
2024 & Educación Parvularia & Seminario Temático de Investigación & 2 & 1.75\\
\addlinespace
2024 & Psicología & Seminario de Práctica 1 Psicología Clínica Cognitivo Constructivista & 1 & 1.00\\
2024 & Psicología & Práctica Profesional 1 Psicología Clínica Cognitivo Constructivista & 1 & 3.00\\
2024 & Sociología & Seminario de Grado & 1 & 1.00\\
2025 & Antropología & Teoría Social & 1 & 1.00\\
2025 & Antropología & Taller de Investigación II - Arqueología & 1 & 1.00\\
\addlinespace
2025 & Antropología & Taller de Investigación II - Social & 1 & 1.00\\
2025 & Antropología & Taller de Memoria II: Mención Social & 1 & 3.00\\
2025 & Antropología & Taller de Memoria I Mención Arqueología & 2 & 3.20\\
2025 & Antropología & Taller de Memoria II: Mención Arqueología & 1 & 3.50\\
2025 & Antropología & Taller de Investigación II - Arqueología & 2 & 3.50\\
\addlinespace
2025 & Antropología & Taller de Investigación I - Social & 2 & 3.65\\
2025 & Antropología & Taller de Investigación II - Arqueología & 6 & 3.72\\
2025 & Psicología & Filosofía & 9 & 3.32\\
2025 & Psicología & Epistemología de las Ciencias Sociales & 2 & 3.40\\
2025 & Psicología & Estadística 1 & 2 & 3.55\\
\addlinespace
2025 & Sociología & Historia Social de América Latina & 5 & 3.60\\
2025 & Sociología & Teorías Sociológicas Contemporáneas & 43 & 3.91\\
2025 & Sociología & Estrategias de Investigación Cuantitativa & 42 & 3.95\\
2025 & Trabajo Social & Estadística I & 2 & 2.70\\
2025 & Trabajo Social & Estado y Sociedad Contemporánea & 2 & 3.10\\
\addlinespace
2025 & Trabajo Social & Estadística I & 9 & 3.56\\
2025 & Trabajo Social & Enfoques Críticos en Trabajo Social & 1 & 3.70\\
2025 & Trabajo Social & Teoría Sociológica II & 3 & 3.73\\
\bottomrule

\end{longtable}

\begin{longtable}[t]{lllr}

\caption{\label{tbl-conteo-secciones}Listado de cursos con más de dos
secciones}

\tabularnewline

\toprule
\cellcolor[HTML]{952626}{\textcolor{white}{\textbf{Año}}} & \cellcolor[HTML]{952626}{\textcolor{white}{\textbf{Carrera}}} & \cellcolor[HTML]{952626}{\textcolor{white}{\textbf{Curso}}} & \cellcolor[HTML]{952626}{\textcolor{white}{\textbf{N secciones}}}\\
\midrule
2018 & Sociología & Introducción a la Sociología & 3\\
2020 & Sociología & Introducción a la Sociología & 3\\
2020 & Sociología & Sociología de las Políticas Públicas & 3\\
2020 & Sociología & Teoría Sociológica Clásica & 3\\
2021 & Antropología & Taller de Investigación I - Social & 12\\
\addlinespace
2021 & Antropología & Taller de Investigación II - Arqueología & 12\\
2021 & Antropología & Taller de Investigación II - Social & 12\\
2021 & Educación Parvularia & Seminario Temático de Investigación & 19\\
2021 & Sociología & Seminario de Grado & 26\\
2021 & Sociología & Transformaciones Sociales del Chile Contemporáneo & 4\\
\addlinespace
2021 & Sociología & Diseños de Investigación & 3\\
2021 & Sociología & Introducción a la Sociología & 3\\
2021 & Sociología & Sociología de las Políticas Públicas & 3\\
2021 & Sociología & Teoría Sociológica Clásica & 3\\
2022 & Antropología & Taller de Investigación I - Social & 16\\
\addlinespace
2022 & Antropología & Taller de Investigación II - Social & 13\\
2022 & Antropología & Taller de Memoria I: Mención Social & 13\\
2022 & Antropología & Taller de Investigación II - Arqueología & 12\\
2022 & Antropología & Taller de Memoria I Mención Arqueología & 11\\
2022 & Antropología & Taller de Memoria II: Mención Arqueología & 11\\
\addlinespace
2022 & Antropología & Taller de Memoria II: Mención Social & 11\\
2022 & Antropología & Taller de Memoria I: Mención Física & 5\\
2022 & Antropología & Taller de Memoria II: Mención Física & 4\\
2022 & Antropología & Taller de Investigación II - Física & 3\\
2022 & Educación Parvularia & Seminario Temático de Investigación & 20\\
\addlinespace
2022 & Educación Parvularia & Seminario Título I & 16\\
2022 & Educación Parvularia & Seminario Título II & 15\\
2022 & Sociología & Seminario de Grado & 27\\
2022 & Sociología & Seminario de Título I & 26\\
2022 & Sociología & Seminario de Título II & 26\\
\addlinespace
2022 & Sociología & Práctica Profesional & 25\\
2022 & Sociología & Transformaciones Sociales del Chile Contemporáneo & 4\\
2022 & Sociología & Introducción a la Sociología & 3\\
2022 & Sociología & Teoría Sociológica Clásica & 3\\
2022 & Trabajo Social & Núcleo Temático Electivo I + D I: & 5\\
\addlinespace
2022 & Trabajo Social & Núcleo Temático Electivo I + D II: & 5\\
2022 & Trabajo Social & Núcleo Temático Electivo I + D III & 5\\
2022 & Trabajo Social & Núcleo Temático Electivo I+D IV & 5\\
2022 & Trabajo Social & Núcleo de Titulación I+D I: & 5\\
2022 & Trabajo Social & Núcleo de Titulación I+D II: & 5\\
\addlinespace
2023 & Antropología & Taller de Investigación II - Social & 18\\
2023 & Antropología & Taller de Investigación I - Social & 16\\
2023 & Antropología & Taller de Investigación II - Arqueología & 15\\
2023 & Antropología & Taller de Memoria II: Mención Social & 14\\
2023 & Antropología & Taller de Memoria I: Mención Social & 13\\
\addlinespace
2023 & Antropología & Taller de Memoria I Mención Arqueología & 11\\
2023 & Antropología & Taller de Memoria II: Mención Arqueología & 11\\
2023 & Antropología & Taller de Investigación II - Física & 4\\
2023 & Antropología & Taller de Memoria I: Mención Física & 3\\
2023 & Antropología & Taller de Memoria II: Mención Física & 3\\
\addlinespace
2023 & Educación Parvularia & Seminario Temático de Investigación & 19\\
2023 & Educación Parvularia & Seminario Título II & 16\\
2023 & Educación Parvularia & Seminario Título I & 14\\
2023 & Sociología & Seminario de Título I & 28\\
2023 & Sociología & Seminario de Título II & 28\\
\addlinespace
2023 & Sociología & Seminario de Grado & 26\\
2023 & Sociología & Práctica Profesional & 22\\
2023 & Sociología & Transformaciones Sociales del Chile Contemporáneo & 4\\
2023 & Trabajo Social & Núcleo Temático Electivo I + D I: & 5\\
2023 & Trabajo Social & Núcleo Temático Electivo I + D II: & 5\\
\addlinespace
2023 & Trabajo Social & Núcleo Temático Electivo I + D III & 5\\
2023 & Trabajo Social & Núcleo Temático Electivo I+D IV & 5\\
2023 & Trabajo Social & Núcleo de Titulación I+D I: & 5\\
2023 & Trabajo Social & Núcleo de Titulación I+D II: & 5\\
2024 & Antropología & Taller de Memoria I: Mención Social & 15\\
\addlinespace
2024 & Antropología & Taller de Investigación II - Arqueología & 14\\
2024 & Antropología & Taller de Investigación II - Social & 14\\
2024 & Antropología & Taller de Investigación I - Social & 13\\
2024 & Antropología & Taller de Memoria II: Mención Social & 13\\
2024 & Antropología & Taller de Memoria II: Mención Arqueología & 11\\
\addlinespace
2024 & Antropología & Taller de Memoria I Mención Arqueología & 10\\
2024 & Antropología & Taller de Investigación II - Física & 8\\
2024 & Antropología & Taller de Memoria I: Mención Física & 3\\
2024 & Antropología & Taller de Memoria II: Mención Física & 3\\
2024 & Educación Parvularia & Seminario Temático de Investigación & 22\\
\addlinespace
2024 & Educación Parvularia & Seminario Título II & 14\\
2024 & Educación Parvularia & Seminario Título I & 10\\
2024 & Educación Parvularia & Práctica Profesional 1 Ciclo & 9\\
2024 & Educación Parvularia & Taller de Práctica 1 & 9\\
2024 & Educación Parvularia & Práctica Profesional 2 Ciclo & 8\\
\addlinespace
2024 & Educación Parvularia & Taller de Práctica 2 & 6\\
2024 & Psicología & Práctica Profesional & 91\\
2024 & Psicología & Memoria de Título & 74\\
2024 & Psicología & Práctica Profesional 2 Psicología Jurídica & 6\\
2024 & Psicología & Práctica Profesional Psicología Jurídica & 6\\
\addlinespace
2024 & Psicología & Seminario de Práctica 2 Psicología Jurídica & 6\\
2024 & Psicología & Seminario de Práctica Psicología Clínica Sistémico & 6\\
2024 & Psicología & Seminario de Práctica Psicología Jurídica & 6\\
2024 & Psicología & Práctica Profesional 2 Psicología Clínica Sistémica & 5\\
2024 & Psicología & Práctica Profesional Psicología Clínica Sistémica & 5\\
\addlinespace
2024 & Psicología & Seminario de Práctica 2 Psicología Clínica Sistémico & 5\\
2024 & Psicología & Práctica Profesional 2 Psicología Educacional & 4\\
2024 & Psicología & Práctica Profesional Psicología Educacional & 4\\
2024 & Psicología & Seminario de Práctica 1 Psicología Educacional & 4\\
2024 & Psicología & Seminario de Práctica 2 Psicología Educacional & 4\\
\addlinespace
2024 & Psicología & Práctica Profesional 2 Psicología del Trabajo y las Organizaciones & 3\\
2024 & Psicología & Práctica Profesional Psicología del Trabajo y las Organizaciones & 3\\
2024 & Psicología & Seminario de Práctica 1 Psicología del Trabajo y las Organizaciones & 3\\
2024 & Psicología & Seminario de Práctica 2 Psicología del Trabajo y las Organizaciones & 3\\
2024 & Sociología & Seminario de Grado & 28\\
\addlinespace
2024 & Sociología & Seminario de Título II & 26\\
2024 & Sociología & Seminario de Título I & 25\\
2024 & Sociología & Práctica Profesional & 24\\
2024 & Sociología & Transformaciones Sociales del Chile Contemporáneo & 4\\
2024 & Trabajo Social & NÚCLEO TEMÁTICO ELECTIVO I+D IV & 11\\
\addlinespace
2024 & Trabajo Social & NUCLEO DE TITULACIÓN I+D II & 9\\
2024 & Trabajo Social & NÚCLEO TEMÁTICO ELECTIVO I+D II: & 7\\
2024 & Trabajo Social & Núcleo Temático Electivo I + D II: & 6\\
2024 & Trabajo Social & Núcleo Temático Electivo I + D I: & 5\\
2024 & Trabajo Social & Núcleo Temático Electivo I + D III & 5\\
\addlinespace
2024 & Trabajo Social & Núcleo de Titulación I+D I: & 5\\
2024 & Trabajo Social & NÚCLEO TEMÁTICO ELECTIVO I+D IV: & 3\\
2024 & Trabajo Social & Núcleo de Titulación I+D II: & 3\\
2025 & Antropología & Taller de Investigación I - Social & 16\\
2025 & Antropología & Taller de Investigación II - Social & 15\\
\addlinespace
2025 & Antropología & Taller de Memoria I: Mención Social & 15\\
2025 & Antropología & Taller de Investigación II - Arqueología & 14\\
2025 & Antropología & Taller de Memoria I Mención Arqueología & 13\\
2025 & Antropología & Taller de Memoria II: Mención Arqueología & 12\\
2025 & Antropología & Taller de Memoria II: Mención Social & 12\\
\addlinespace
2025 & Antropología & Taller de Investigación II - Física & 7\\
2025 & Antropología & Taller de Memoria II: Mención Física & 7\\
2025 & Antropología & Taller de Memoria I: Mención Física & 6\\
2025 & Educación Parvularia & Seminario Temático de Investigación & 24\\
2025 & Educación Parvularia & Seminario Título I & 21\\
\addlinespace
2025 & Educación Parvularia & Seminario Título II & 16\\
2025 & Educación Parvularia & Práctica Profesional 1 Ciclo & 10\\
2025 & Educación Parvularia & Taller de Práctica 1 & 10\\
2025 & Educación Parvularia & Práctica Profesional 2 Ciclo & 9\\
2025 & Educación Parvularia & Taller de Práctica 2 & 9\\
\addlinespace
2025 & Psicología & Memoria de Título & 16\\
2025 & Psicología & Práctica Profesional & 16\\
2025 & Psicología & Práctica Profesional 1 Psicología Clínica Psicoanálisis & 6\\
2025 & Psicología & Práctica Profesional 1 Psicología Jurídica & 6\\
2025 & Psicología & Práctica Profesional 2 Psicología Clínica Psicoanálisis & 6\\
\addlinespace
2025 & Psicología & Práctica Profesional 2 Psicología Jurídica & 6\\
2025 & Psicología & Seminario de Práctica 1 Psicología Clínica Psicoanálisis & 6\\
2025 & Psicología & Seminario de Práctica 1 Psicología Jurídica & 6\\
2025 & Psicología & Seminario de Práctica 2 Psicología Clínica Psicoanálisis & 6\\
2025 & Psicología & Seminario de Práctica 2 Psicología Jurídica & 6\\
\addlinespace
2025 & Psicología & Práctica Profesional 1 Psicología Clínica Sistémica & 5\\
2025 & Psicología & Práctica Profesional 2 Psicología Clínica Sistémica & 5\\
2025 & Psicología & Seminario de Práctica 1 Psicología Clínica Sistémico & 5\\
2025 & Psicología & Seminario de Práctica 2 Psicología Clínica Sistémico & 5\\
2025 & Psicología & Práctica Profesional 1 Psicología Educacional & 4\\
\addlinespace
2025 & Psicología & Práctica Profesional 2 Psicología Educacional & 4\\
2025 & Psicología & Seminario de Práctica 1 Psicología Educacional & 4\\
2025 & Psicología & Seminario de Práctica 2 Psicología Educacional & 4\\
2025 & Psicología & Práctica Profesional 2 Psicología del Trabajo y las Organizaciones & 3\\
2025 & Psicología & Seminario de Práctica 2 Psicología del Trabajo y las Organizaciones & 3\\
\addlinespace
2025 & Sociología & Seminario de Grado & 31\\
2025 & Sociología & Seminario de Título I & 27\\
2025 & Sociología & Seminario de Título II & 26\\
2025 & Sociología & Práctica Profesional & 25\\
2025 & Sociología & Transformaciones Sociales del Chile Contemporáneo & 3\\
\addlinespace
2025 & Trabajo Social & NUCLEO DE TITULACIÓN I+D II & 18\\
2025 & Trabajo Social & Núcleo de Titulación I+D I: & 18\\
2025 & Trabajo Social & NÚCLEO TEMÁTICO ELECTIVO I+D II: & 16\\
2025 & Trabajo Social & NÚCLEO TEMÁTICO ELECTIVO I+D IV & 16\\
2025 & Trabajo Social & Núcleo Temático Electivo I + D I: & 16\\
\addlinespace
2025 & Trabajo Social & Núcleo Temático Electivo I + D III & 16\\
\cellcolor[HTML]{f2f2f2}{\textbf{2026}} & \cellcolor[HTML]{f2f2f2}{\textbf{Sociología}} & \cellcolor[HTML]{f2f2f2}{\textbf{Práctica Profesional}} & \cellcolor[HTML]{f2f2f2}{\textbf{3}}\\
\bottomrule

\end{longtable}


\backmatter

% --- Cuerpo del libro (capítulos) en arábigos ---
\mainmatter
\pagestyle{scrheadings}

% --- (Opcional) Estilo del índice general ---
% \addtocontents{toc}{\protect\thispagestyle{plain}}


\end{document}
